\chapter{运动规划导论}
\label{ch:planning}

% ════════════════════════════════════════════════════════════════
%  规划导论：全吸收章。
%  吸收 LaValle《Planning Algorithms》2006（Ch2 离散规划 / Ch4 构型空间 / Ch5 采样规划）
%  + Karaman-Frazzoli 2011（PRM*/RRG/RRT* 渐近最优）+ Hart-Nilsson-Raphael 1968（A*）
%  + Kuffner-LaValle 2000（RRT-Connect）+ Stentz 1994（D*）+ Zucker 2013（CHOMP）
%  + Solovey 2020（RRT* 最优性勘误）。自包含独立记录，读者无需翻原书。
%  本章主要据权威教材/论文重建（个人笔记未同步），存疑处打 \rebuilt。
%  教学层遵循 v5.0 理论教学规范。
% ════════════════════════════════════════════════════════════════

\begin{practice}[前置自测]
开始之前，先试答下列问题。能流畅作答者可只速读本章表格与速查卡；卡住之处，正是本章要为你解决的。
\begin{enumerate}
  \item 一个在平面上既能平移又能旋转的刚体机器人，其“所有可能位姿”构成的集合是什么空间？它和 $\SE(2)$ 有何关系？
  \item 为什么不能简单地把整个工作空间划成细栅格、再对每个格点跑一遍搜索，就解决六自由度机械臂的避障规划？
  \item Dijkstra 与 A$^*$ 的唯一实质区别是什么？A$^*$ 的启发函数 $h$ 满足什么条件才能保证找到最优路径？
  \item RRT（快速探索随机树）凭什么能“快速”铺满空间？它找到的路径是最优的吗？
  \item 给定一条已经无碰撞但很“锯齿”的路径，如何把它优化成光滑、可执行的轨迹？这属于哪一类规划方法？
\end{enumerate}
\end{practice}

\section*{本章目标}
学完本章，你应当能够：
\begin{enumerate}
  \item \textbf{建立}构型空间 $\mathcal{C}$、自由空间 $\mathcal{C}_{free}$、障碍空间 $\mathcal{C}_{obs}$ 的概念，并推导 $\SO(2)\cong S^1$、$\SO(3)\cong\mathbb{RP}^3$、$\SE(3)=\mathbb{R}^3\times\mathbb{RP}^3$ 的拓扑结构；
  \item \textbf{陈述}钢琴搬运工问题（基本运动规划问题）及其 PSPACE-hard 复杂度，说明它为何驱动了采样式方法；
  \item \textbf{推导}图搜索的统一前向框架，并区分 BFS/DFS/Dijkstra/A$^*$ 仅在优先队列排序函数上的不同；
  \item \textbf{证明} Dijkstra 出队即最优、A$^*$ 在可采纳启发下的最优性、一致启发蕴含可采纳且 $f$ 沿路径单调；
  \item \textbf{写出} RRT、RRT-Connect、PRM、RRT$^*$、RRG、PRM$^*$ 的完整算法，并解释 Voronoi 偏置与“选最优父 + 重接线”两步；
  \item \textbf{判别}各采样算法的概率完备性与渐近最优性，掌握临界连接半径 $r(n)\sim(\log n/n)^{1/d}$ 的几何来由；
  \item \textbf{写出}轨迹优化的一般泛函形式，推导 CHOMP 的协变梯度更新，并说明“全局采样 + 局部优化”的级联流水线。
\end{enumerate}

\subsection*{本章知识导航}
本章是一条“\emph{把连续避障问题逐级离散化、再在离散结构上搜索或优化}”的主线。三大范式（搜索式 / 采样式 / 优化式）本质上是\textbf{对 $\mathcal{C}_{free}$ 离散化粒度不同}的同一框架。结构如下：

\begin{center}
\small
\begin{tikzpicture}[
  node distance=5mm and 7mm, >=Stealth,
  blk/.style={draw, rounded corners, align=center, font=\small, inner sep=3pt, fill=main!6, draw=main!60!black},
  grp/.style={draw, rounded corners, align=center, font=\small, inner sep=3pt, fill=second!6, draw=second!60!black},
  opt/.style={draw, rounded corners, align=center, font=\small, inner sep=3pt, fill=third!8, draw=third!60!black},
  ln/.style={->, thick, gray!60!black}]
\node[blk] (cspace) {构型空间\\ $\mathcal{C},\mathcal{C}_{free},\mathcal{C}_{obs}$};
\node[blk, right=of cspace] (disc) {离散化\\（栅格 / 采样）};
\node[grp, right=of disc] (search) {搜索式\\ Dijkstra, A$^*$, D$^*$};
\node[grp, below=7mm of search] (sample) {采样式\\ PRM, RRT, RRT$^*$};
\node[opt, below=7mm of sample] (traj) {优化式\\ CHOMP, TrajOpt};
\node[blk, below=7mm of cspace, fill=second!6, draw=second!60!black] (dp) {动态规划\\ Bellman 递推};
\draw[ln] (cspace)--(disc);
\draw[ln] (disc)--(search);
\draw[ln] (disc.south) |- (sample.west);
\draw[ln] (dp.east) -- (search.west|-dp.east) ;
\draw[ln] (dp)--(cspace);
\draw[ln] (sample) -- (traj);
\draw[ln] (search.south) -- (sample.north);
\end{tikzpicture}
\end{center}

\noindent\textbf{推荐路径}：第 \ref{sec:plan-cspace} 节（构型空间）是一切的地基；第 \ref{sec:plan-search}--\ref{sec:plan-astar} 节（图搜索 + A$^*$，含最优性证明）是\emph{搜索式}主干，也是采样式查询阶段复用的工具；第 \ref{sec:plan-grid} 节（栅格启发）面向工程；第 \ref{sec:plan-sampling}--\ref{sec:plan-optimal} 节（采样式 + 渐近最优理论）是本章\emph{篇幅最大、最核心}的部分，VIO/导航/操作臂规划都依赖它；第 \ref{sec:plan-traj} 节（轨迹优化）是收尾，衔接后续控制章。带 $*$ 的 D$^*$、渐近最优完整证明为进阶/选读。

\subsection*{前置知识桥接}
本章只用到通用的拓扑、图论、概率与微积分基础，以及本书 \cref{ch:lie} 给出的李群语言。具体地：当构型 $\mathbf{q}$ 含姿态时（如 $\mathcal{C}=\SE(3)$ 的刚体规划），构型空间上的“距离度量”“插值（Steer）”须用 \cref{ch:lie} 的 $\mathrm{Exp}/\mathrm{Log}$ 与测地线实现——本章源文献多在 $\mathbb{R}^d$ 上叙述，移植到 $\SO(3)/\SE(3)$ 时距离、采样需改为流形版（\cref{sec:plan-cspace} 末有专门说明）。\cref{ch:nlopt}（非线性优化）的最小二乘 + 高斯牛顿框架，将在 \cref{sec:plan-traj} 的轨迹优化里再次出现。

\subsection*{如果跳过本章会怎样}
\begin{itemize}
  \item 在做机械臂或移动机器人导航时，你会写出“从 A 点运动到 B 点且避开障碍”的需求，却\emph{不知道}它在数学上等价于“在 $\mathcal{C}_{free}$ 中找一条连续路径”，从而把一个本可统一处理的问题，对每种机器人各写一套特例代码；
  \item 你会直接调用 OMPL、MoveIt 里的 \texttt{RRTConnect}、\texttt{RRTstar} 求解器，却\emph{不理解}为什么 RRT 给出的路径“能用但不漂亮”、为什么 RRT$^*$ “越跑越好”、为什么连接半径要写成 $(\log n/n)^{1/d}$——只能照抄参数而不会调试与选型。
\end{itemize}

\begin{table}[H]\centering\small
\caption{本章阅读方式建议}
\label{tab:plan-reading}
\begin{tabular}{lll}
\toprule
阅读方式 & 时间 & 适合谁 \\
\midrule
精读（含全部推导与证明） & 6--8 小时 & 首次系统学习、要做规划算法实现/研究 \\
速读（跳过渐近最优完整证明） & 2.5 小时 & 有规划基础、想对齐本书记号与三范式脉络 \\
速查（只看小结的符号表 + 定理速查表 + 算法清单） & 20 分钟 & 写代码、选规划器时回来查 \\
\bottomrule
\end{tabular}
\end{table}

\begin{note}
\textbf{本章记号与约定.}\;
连续构型空间记 $\mathcal{C}$（LaValle\cite{lavalle2006planning} 用 $\mathcal{C}$，Karaman--Frazzoli\cite{karaman2011sampling} 用 $X=(0,1)^d$；本书统一 $\mathcal{C}$），其元素\emph{构型} $\mathbf{q}\in\mathcal{C}$；离散状态/图用 $X$ 与图 $G=(V,E)$，状态 $\mathbf{x}\in X$。自由空间 $\mathcal{C}_{free}$、障碍空间 $\mathcal{C}_{obs}=\mathcal{C}\setminus\mathcal{C}_{free}$。起点 $\mathbf{q}_I$、目标 $\mathbf{q}_G$（或目标集 $\mathcal{C}_{goal}$）。维数 $d=\dim\mathcal{C}$。
A$^*$ 沿用 AI 经典记号：已走代价 $g(\mathbf{x})$（cost-to-come）、启发 $h(\mathbf{x})$（cost-to-go 下估计，真值 $h^*$）、评价 $f=g+h$。
采样常数：连接半径 $r(n)$、Steer 步长 $\eta$、单位球体积 $\zeta_d=\pi^{d/2}/\Gamma(d/2+1)$、Lebesgue 测度 $\mu$。
\textbf{字母隔离声明}：本章 $\mu$ 指\emph{测度}（非均值）、$\lambda$ 指\emph{渗流强度/权衡系数}（非特征值或波长）、代价函数 $c$（非 Barfoot 的旋转别名 $\mathbf{C}$，本书旋转一律 $\mathbf{R}$），均仅在本章局部生效，与 SLAM 章的同字母无关。
\end{note}

% ════════════════════════════════════════════════════════════════
\section{构型空间}
\label{sec:plan-cspace}

\paragraph{动机：把“刚体避障”变成“点找路”。}
考虑一个在平面上推钢琴的搬运工：钢琴是一个有形状、会平移会旋转的刚体，要从房间一头挪到另一头、绕过家具。直接在工作空间里推理“钢琴的每个点是否撞到家具”非常繁琐，因为机器人有体积、有朝向。运动规划的核心抽象\cite{lavalle2006planning}是：把机器人\emph{所有可能位姿}的集合本身当作一个空间——\textbf{构型空间}（configuration space，C-space）$\mathcal{C}$——于是机器人在 $\mathcal{C}$ 中退化成\emph{一个点}，“刚体避障”就变成了“在 $\mathcal{C}$ 中绕开障碍区、把一个点连到另一个点”。这个抽象由 Lozano-Pérez 开创、Latombe 的教材统一，威力在于：无论机器人多复杂（多连杆、多自由度），一旦写到 $\mathcal{C}$ 上，规划问题就被统一成同一个数学问题。

\paragraph{反面：不抽象会怎样。}
若不引入 $\mathcal{C}$，你只能针对每种机器人（点、圆盘、长杆、机械臂）各写一套“几何避障”逻辑，且无法回答“是否存在路径”这种全局问题——因为你没有一个统一的“位姿空间”可供搜索。更糟的是，旋转带来的拓扑（绕一圈回到原点）若用欧拉角硬编码，会引入万向锁等假奇异（见下文 $\SO(3)$ 讨论），让规划器在某些姿态附近行为错乱。

\subsection{基本拓扑：开闭集、流形、连通性}
\label{sec:plan-topology}

规划要在连续空间上“把一点连到另一点”，需要拓扑语言来精确描述“连续”“逼近边界”“连通”。以下定义来自 LaValle\cite{lavalle2006planning} §4.1。

\begin{definition}[拓扑空间与开闭集]\label{def:plan-topology}
集合 $X$ 配一族子集（称\textbf{开集}），满足三公理则为\textbf{拓扑空间}：(1) 任意多个（含无穷多）开集之并仍是开集；(2) \emph{有限}个开集之交仍是开集；(3) $X$ 与 $\varnothing$ 都是开集。子集 $C\subseteq X$ 称\textbf{闭集}当且仅当其补 $X\setminus C$ 是开集。
\end{definition}

\noindent 开区间 $(0,1)$ 含端点之间所有实数却不含端点：可构造序列 $\tfrac12,\tfrac14,\dots,\tfrac1{2^i}\to0$ 任意逼近边界却取不到；闭区间 $[0,1]$ 含端点。注意：无穷多开集之\emph{交}未必开（如 $\bigcap_i(-\tfrac1i,\tfrac1i)=\{0\}$ 闭），这正是公理 (2) 只对有限交成立的原因。一个反直觉点：并非每个子集非开即闭——$\mathbb{R}$ 上 $[0,2\pi)$ 既不开也不闭；而任何拓扑空间里 $X$ 与 $\varnothing$ \emph{既开又闭}。

\begin{definition}[特殊点与闭包]\label{def:plan-points}
设 $U\subseteq X$，$\mathbf{x}\in X$：若存在开集 $O_1$ 使 $\mathbf{x}\in O_1\subseteq U$，称 $\mathbf{x}$ 为 $U$ 的\textbf{内点}，全体内点之集为\textbf{内部} $\mathrm{int}(U)$；若存在开集 $O_2$ 使 $\mathbf{x}\in O_2\subseteq X\setminus U$，称\textbf{外点}；既非内点也非外点者为\textbf{边界点}，全体边界点之集为\textbf{边界} $\partial U$。内点与边界点统称\textbf{极限点}，全体极限点之集是闭集，称\textbf{闭包}
\begin{equation}
  \mathrm{cl}(U)=\mathrm{int}(U)\cup\partial U.
  \label{eq:plan-closure}
\end{equation}
\end{definition}

\noindent 由 \cref{def:plan-points}：闭集含其全部边界点（闭包总把边界点加进去）；开集不含任一自身边界点。这一区别后文极关键——$\mathcal{C}_{free}$ 是开集（机器人可任意逼近障碍却不触碰），但“最短路径”这种优化在开集上\emph{无定义}（下确界取不到），须改用闭包 $\mathrm{cl}(\mathcal{C}_{free})$。

\begin{definition}[连续与同胚]\label{def:plan-homeo}
$f:X\to Y$ \textbf{连续}当且仅当对 $Y$ 每个开集 $O$，预像 $f^{-1}(O)=\{\mathbf{x}\in X\mid f(\mathbf{x})\in O\}$ 是 $X$ 的开集。若 $f$ 是双射且 $f$ 与 $f^{-1}$ 都连续，则 $f$ 为\textbf{同胚}，称 $X,Y$ \textbf{同胚}，记 $X\cong Y$。
\end{definition}

\noindent 这一开集预像定义比经典 $\delta$-$\epsilon$ 简洁：连续函数之复合连续，只需“开集的预像的预像仍开”一行论证。同胚是拓扑里最基本的等价（自反、对称、传递）——“甜甜圈与单柄咖啡杯同胚”即指其表面都恰有一个洞。例：任意两开区间同胚（$(0,1)\to(0,5)$ 用 $x\mapsto5x$）；$x\mapsto\tfrac2\pi\tan^{-1}x$ 使 $(-1,1)\cong\mathbb{R}$（有界与无界可同胚）。\textbf{警告}：$[0,1]$ 与 $(0,1)$ \emph{不}同胚（端点作梗）。

\begin{definition}[流形]\label{def:plan-manifold}
拓扑空间 $M\subseteq\mathbb{R}^m$ 是 $n$ 维\textbf{流形}，若对每个 $\mathbf{x}\in M$ 存在开集 $O\subset M$ 使：(1) $\mathbf{x}\in O$；(2) $O$ 同胚于 $\mathbb{R}^n$；(3) $n$ 对所有 $\mathbf{x}\in M$ 固定。该固定 $n$ 称流形\textbf{维数} $\dim M=n$。
\end{definition}

\noindent 第 (2) 条最关键：流形的定义性质是\emph{局部像欧氏空间}。这与 \cref{ch:lie} 把李群当流形、在单位元附近用切空间做线性运算是同一思想。Whitney 嵌入定理保证任何 $n$ 维流形可嵌入 $\mathbb{R}^{2n+1}$，故取 $M\subseteq\mathbb{R}^m$ 不失一般性。常用低维流形：圆 $S^1=\{(x,y)\mid x^2+y^2=1\}$（与 $(0,1)$ \emph{不}同胚）；$n$ 维球面 $S^n=\{\mathbf{x}\in\mathbb{R}^{n+1}\mid\|\mathbf{x}\|=1\}$；$n$ 维环面 $T^n=(S^1)^n$（注意 $S^1\times S^1\ne S^2$）。

\begin{insight}
$S^1$ 可由区间 $[0,1]$ 作\textbf{等同}（identification）$0\sim1$ 得到（把两端“粘”成闭环），同胚由 $\theta\mapsto(\cos2\pi\theta,\sin2\pi\theta)$ 给出。这种“商拓扑”是描述旋转的关键——旋转 $2\pi$ 回到原点，正是 $0\sim1$ 的环绕。\textbf{规划算法必须识别这种环绕}：在柱面 $\mathcal{C}=\mathbb{R}\times S^1$ 上若不知道“上下边界等同”，就可能误判 $\mathbf{q}_I$ 到 $\mathbf{q}_G$ 不连通。
\end{insight}

\paragraph{实射影空间 $\mathbb{RP}^n$。}
$\mathbb{RP}^n$ 定义为 $\mathbb{R}^{n+1}$ 中过原点的所有直线之集。每条这样的直线与 $S^n$ 恰交于两个\textbf{对径点}（antipodal）。把 $S^n$ 上所有对径点等同，得
\begin{equation}
  \mathbb{RP}^n\cong S^n/\!\sim\quad(\mathbf{x}\sim-\mathbf{x}).
  \label{eq:plan-rpn}
\end{equation}
$\mathbb{RP}^3$ 是运动规划最重要的流形之一——它\emph{就是} $\SO(3)$（见 \cref{sec:plan-so3}）。

\paragraph{连通性与同伦。}
$X$ \textbf{路径连通}：对任意 $\mathbf{x}_1,\mathbf{x}_2\in X$ 存在连续路径 $\tau:[0,1]\to X$ 使 $\tau(0)=\mathbf{x}_1,\tau(1)=\mathbf{x}_2$。两路径 $\tau_1,\tau_2$（端点固定）\textbf{同伦}，若存在连续 $h:[0,1]\times[0,1]\to X$ 使 $h(s,0)=\tau_1(s)$、$h(s,1)=\tau_2(s)$、$h(0,t)$ 与 $h(1,t)$ 恒为端点。$t$ 像旋钮，连续地把 $\tau_1$ 形变到 $\tau_2$，形变中路径\emph{过不去洞}。若 $X$ 所有路径同属一个同伦类则 $X$ \textbf{单连通}，否则\textbf{多连通}。同伦类的概念后文会再现：RRT$^*$ 在含障碍场景里“发现不同同伦类的路径”正是其大幅降低代价的来由（\cref{sec:plan-optimal}）。

\subsection{典型构型空间：$\SO(2),\SE(2),\SO(3),\SE(3)$}
\label{sec:plan-so3}

\paragraph{平面旋转 $\SO(2)\cong S^1$。}
从 $2\times2$ 实矩阵 $\mathbb{R}^4$ 出发，加正交约束 $\mathbf{A}\mathbf{A}^\top=\mathbf{I}$，展开得 $a^2+b^2=1$、$c^2+d^2=1$、$ac+bd=0$，即 1 个列正交方程 + 2 个单位列方程，从 4 维减 3 维得 1 维流形。对每个 $(a,b)\in S^1$ 有两组 $(c,d)=\pm(b,-a)$，是两个不连通的圆 $S^1\sqcup S^1$（即 $O(2)$）；再加 $\det\mathbf{A}=1$ 选其一，得
\begin{equation}
  \SO(2)\cong S^1.
  \label{eq:plan-so2}
\end{equation}
等价地，单位复数 $a+b\mathrm{i}$（$a^2+b^2=1$）在乘法下与 $\SO(2)$ 同构：两旋转复合 = 复数相乘 $e^{\mathrm{i}\theta_1}e^{\mathrm{i}\theta_2}=e^{\mathrm{i}(\theta_1+\theta_2)}$。平面刚体可平移可旋转，故\textbf{平面刚体的 C-space}
\begin{equation}
  \mathcal{C}=\mathbb{R}^2\times S^1=\SE(2)\quad(\dim=3).
  \label{eq:plan-se2}
\end{equation}

\paragraph{空间旋转 $\SO(3)\cong\mathbb{RP}^3$。}
为何不用欧拉角（yaw-pitch-roll $\alpha,\beta,\gamma$）描述 $\SO(3)$？因为它\emph{破坏拓扑}：存在非零角给出单位旋转、存在连续多组角给同一旋转矩阵（万向锁），造成理论与实践困难。正确做法用单位四元数。回顾 \cref{ch:lie} 与 \cref{ch:rigid_body}：四元数 $h=a+b\mathrm{i}+c\mathrm{j}+d\mathrm{k}$（本书 Hamilton 约定，标量 $a=w$ 在前，与 LaValle 一致），满足 $\mathrm{i}^2=\mathrm{j}^2=\mathrm{k}^2=\mathrm{i}\mathrm{j}\mathrm{k}=-1$。单位四元数（$a^2+b^2+c^2+d^2=1$，构成 $S^3$）映到旋转矩阵
\begin{equation}
  \mathbf{R}(h)=\begin{bmatrix}2(a^2+b^2)-1&2(bc-ad)&2(bd+ac)\\2(bc+ad)&2(a^2+c^2)-1&2(cd-ab)\\2(bd-ac)&2(cd+ab)&2(a^2+d^2)-1\end{bmatrix},
  \label{eq:plan-quat2rot}
\end{equation}
绕单位轴 $\mathbf{v}$ 转 $\theta$ 对应 $h=\cos\tfrac\theta2+(\mathbf{v}\sin\tfrac\theta2)$（虚部）。关键观察：$\mathbf{R}(h)=\mathbf{R}(-h)$（绕 $\mathbf{v}$ 转 $\theta$ 等于绕 $-\mathbf{v}$ 转 $2\pi-\theta$，验证 $\cos\tfrac{2\pi-\theta}2=-\cos\tfrac\theta2=-a$、虚部同号取反），故单位四元数到 $\SO(3)$ 是\textbf{二对一}。把 $S^3$ 上对径点等同（$h\sim-h$），由 \cref{eq:plan-rpn} 得
\begin{equation}
  \SO(3)\cong\mathbb{RP}^3=S^3/\!\sim.
  \label{eq:plan-so3rp3}
\end{equation}
因此\textbf{空间刚体的 C-space}（3 平移 + 3 旋转）
\begin{equation}
  \mathcal{C}=\mathbb{R}^3\times\mathbb{RP}^3=\SE(3)\quad(\dim=6),
  \label{eq:plan-se3}
\end{equation}
恰是自由飞行刚体的 6 自由度。

\begin{pitfall}[概念误区：把“正交矩阵”等同于“旋转矩阵”，漏掉 $\det=+1$]
新手常以为 $\mathbf{R}^\top\mathbf{R}=\mathbf{I}$ 就是旋转。实际上 $\mathbf{R}^\top\mathbf{R}=\mathbf{I}$ 只保证 $\mathbf{R}\in O(n)$（正交群），它含 $\det=+1$（旋转 $\SO(n)$）与 $\det=-1$（含反射）两支——这正是上面 $O(2)=S^1\sqcup S^1$ 两个圆、$O(3)=S^3\sqcup S^3$ 两个球面的来由。漏掉 $\det$ 约束会把镜像当旋转，在 $\SO(3)$ 上做指数映射/插值时混入反射会导致不连续、不可积。凡声称“旋转”，必须同时写 $\mathbf{R}^\top\mathbf{R}=\mathbf{I}$ 且 $\det\mathbf{R}=+1$。
\end{pitfall}

\paragraph{多连杆与机械臂。}
$n$ 个自由漂浮刚体的 C-space 是各自的笛卡尔积 $\mathcal{C}=\mathcal{C}_1\times\cdots\times\mathcal{C}_n$。运动学链/树须逐情形分析：$\mathcal{W}=\mathbb{R}^2$ 中 $n$ 个旋转关节、各关节满程 $\theta_i\in[0,2\pi)$ 时 $\mathcal{C}=S^1\times\cdots\times S^1=T^n$（$n$ 维环面）；若各关节受限 $\theta_i\in(-\tfrac\pi2,\tfrac\pi2)$ 则 $\mathcal{C}=\mathbb{R}^n$。\textbf{陷阱}：用 DH 参数化便于变换计算却忽略拓扑——若用三个零长连杆模拟球腕，会重蹈欧拉角覆辙。

\begin{table}[H]\centering\small
\caption{典型机器人的构型空间}
\label{tab:plan-cspaces}
\begin{tabular}{lllc}
\toprule
机器人 & 工作空间 $\mathcal{W}$ & 构型 $\mathbf{q}$ & $\mathcal{C}$（$\dim$） \\
\midrule
平面平移点 & $\mathbb{R}^2$ & $(x,y)$ & $\mathbb{R}^2$ (2) \\
平面刚体 & $\mathbb{R}^2$ & $(x,y,\theta)$ & $\mathbb{R}^2\times S^1=\SE(2)$ (3) \\
空间刚体 & $\mathbb{R}^3$ & $(x,y,z,h)$，$h$ 单位四元数 & $\mathbb{R}^3\times\mathbb{RP}^3=\SE(3)$ (6) \\
$n$ 旋转关节臂 & --- & $(\theta_1,\dots,\theta_n)$ & $T^n=(S^1)^n$ ($n$) \\
\bottomrule
\end{tabular}
\end{table}

\subsection{障碍空间、自由空间与基本规划问题}
\label{sec:plan-cobs}

设工作空间 $\mathcal{W}=\mathbb{R}^2$ 或 $\mathbb{R}^3$ 含\textbf{障碍区} $\mathcal{O}\subset\mathcal{W}$，机器人 $\mathcal{A}\subset\mathcal{W}$（$\mathcal{A},\mathcal{O}$ 用半代数模型，含多边形/多面体）。记 $\mathcal{A}(\mathbf{q})$ 为按构型 $\mathbf{q}$ 变换后的机器人。

\begin{definition}[C-空间障碍区与自由空间]\label{def:plan-cobs}
\begin{equation}
  \mathcal{C}_{obs}=\{\mathbf{q}\in\mathcal{C}\mid\mathcal{A}(\mathbf{q})\cap\mathcal{O}\ne\varnothing\},\qquad
  \mathcal{C}_{free}=\mathcal{C}\setminus\mathcal{C}_{obs}.
  \label{eq:plan-cobsfree}
\end{equation}
即 $\mathcal{C}_{obs}$ 是所有使变换后机器人与障碍相交的构型之集。因 $\mathcal{O}$ 与 $\mathcal{A}(\mathbf{q})$ 在 $\mathcal{W}$ 中是\emph{闭集}，$\mathcal{C}_{obs}$ 是 $\mathcal{C}$ 中\textbf{闭集}，故 $\mathcal{C}_{free}$ 是\textbf{开集}。
\end{definition}

\noindent “接触”边界条件：$\mathbf{q}\in\mathcal{C}_{obs}$ 当 $\mathrm{int}(\mathcal{O})\cap\mathrm{int}(\mathcal{A}(\mathbf{q}))=\varnothing$ 但 $\mathcal{O}\cap\mathcal{A}(\mathbf{q})\ne\varnothing$（仅边界相交也算碰撞）。对 $m$ 连杆机器人（含自碰撞），设\emph{碰撞对集} $P$（$(i,j)\in P$ 表不许 $\mathcal{A}_i,\mathcal{A}_j$ 相交）：
\begin{equation}
  \mathcal{C}_{obs}=\Big(\bigcup_{i=1}^m\{\mathbf{q}\mid\mathcal{A}_i(\mathbf{q})\cap\mathcal{O}\ne\varnothing\}\Big)\cup\Big(\bigcup_{[i,j]\in P}\{\mathbf{q}\mid\mathcal{A}_i(\mathbf{q})\cap\mathcal{A}_j(\mathbf{q})\ne\varnothing\}\Big).
  \label{eq:plan-cobs-multi}
\end{equation}
$|P|=O(m^2)$，但几何分析常可剔除大量不可能碰撞的对（相邻连杆恒接触者不入 $P$）。

\begin{definition}[基本运动规划问题（钢琴搬运工问题）]\label{def:plan-piano}
给定 (1) 工作空间 $\mathcal{W}=\mathbb{R}^2$ 或 $\mathbb{R}^3$；(2) 半代数障碍区 $\mathcal{O}\subset\mathcal{W}$；(3) 半代数机器人（刚体 $\mathcal{A}$ 或多连杆 $\mathcal{A}_1,\dots,\mathcal{A}_m$）；(4) 由“可施加的全部变换”导出的 $\mathcal{C}$ 及 $\mathcal{C}_{obs},\mathcal{C}_{free}$；(5) 初始构型 $\mathbf{q}_I\in\mathcal{C}_{free}$；(6) 目标构型 $\mathbf{q}_G\in\mathcal{C}_{free}$（$(\mathbf{q}_I,\mathbf{q}_G)$ 称\emph{查询对}）。\textbf{任务}：计算连续路径 $\tau:[0,1]\to\mathcal{C}_{free}$ 使 $\tau(0)=\mathbf{q}_I,\tau(1)=\mathbf{q}_G$，或正确报告不存在。
\end{definition}

\begin{theorem}[复杂度：PSPACE-hard]\label{thm:plan-pspace}
基本运动规划问题（广义钢琴搬运工问题）是 \textbf{PSPACE-hard}（Reif 1979\cite{lavalle2006planning}），蕴含 NP-hard。难点在于 $\dim\mathcal{C}$ 无界——完备的精确几何算法存在（Schwartz--Sharir、Canny），但其复杂度随维数爆炸，不适用于实践。
\end{theorem}

\begin{insight}
\cref{thm:plan-pspace} 是采样式规划存在的根本动因。显式构造 $\mathcal{C}_{obs}$ 或 $\mathcal{C}_{free}$ 的边界既不直接、又随维数指数爆炸。但有一个不对称的好消息：\emph{显式算出 $\mathcal{C}_{free}$ 极难，而测试某个给定构型 $\mathbf{q}$ 是否属于 $\mathcal{C}_{free}$ 却很高效}（一次碰撞检测）。整个采样式规划就立足于这一不对称——把碰撞检测当黑箱，只问“此构型/此线段是否无碰”，绝不显式建 $\mathcal{C}_{obs}$。
\end{insight}

\paragraph{纯平移情形可显式构造（Minkowski 差）。}
作为对照，当机器人只平移（$\mathcal{C}=\mathbb{R}^n$）时 $\mathcal{C}_{obs}$ 有闭式：$\mathcal{C}_{obs}=\mathcal{O}\ominus\mathcal{A}(\mathbf{0})=\mathcal{O}\oplus(-\mathcal{A}(\mathbf{0}))$，其中 Minkowski 和 $X\oplus Y=\{\mathbf{x}+\mathbf{y}\}$，$-\mathcal{A}$ 把 $\mathcal{A}$ 每点取负。例（一维）：$\mathcal{W}=\mathbb{R}$，机器人 $\mathcal{A}=[-1,2]$、障碍 $\mathcal{O}=[0,4]$，$-\mathcal{A}=[-2,1]$，则 $\mathcal{C}_{obs}=\mathcal{O}\oplus(-\mathcal{A})=[-2,5]$。凸多边形机器人 + 凸多边形障碍时 $\mathcal{C}_{obs}$ 仍凸，可由“按内/外法向角在 $S^1$ 上归并”的 $O(n+m)$ 算法精确求出。一旦加入旋转（须在 $\mathbb{R}^2\times S^1$ 中按 $\theta$ 切片，每片是平移情形），代数迅速复杂化——再次印证高维显式 $\mathcal{C}_{obs}$ 不可行。

\paragraph{移植到流形的接口。}\rebuilt
本章源文献多在 $\mathbb{R}^d$ 上叙述。当 $\mathcal{C}=\SO(3)$ 或 $\SE(3)$ 时，欧氏距离 $\|\mathbf{x}-\mathbf{y}\|$ 须换为李群测地距离（如 $\SO(3)$ 上 $\|\mathrm{Log}(\mathbf{R}_x^\top\mathbf{R}_y)\|$，见 \cref{ch:lie}）；单位球体积 $\zeta_d$、测度 $\mu$ 须换为相应黎曼体积。均匀采样 $\SO(3)$ 须用四元数的 Haar 测度（取 $u_1,u_2,u_3\in[0,1]$ 均匀，$h=(\sqrt{1-u_1}\sin2\pi u_2,\sqrt{1-u_1}\cos2\pi u_2,\sqrt{u_1}\sin2\pi u_3,\sqrt{u_1}\cos2\pi u_3)$）——\emph{欧拉角均匀采样不给 $\SO(3)$ 均匀}。本章的概率完备性/渐近最优性结论在“局部欧氏 + 测度良性”条件下推广成立（源文献明言结论对任意“密度有正下界的绝对连续分布”都成立）。

\begin{practice}
\begin{enumerate}
  \item 由矩阵群法推导 $\SO(2)\cong S^1$（写出列正交 + 单位列方程，数清自由度）。
  \item 验证 $\mathbf{R}(h)=\mathbf{R}(-h)$（用 \cref{eq:plan-quat2rot} 或轴角形式），并解释这为何给出 $\SO(3)\cong\mathbb{RP}^3$ 的二对一。
  \item 平面 $L$ 形机器人在带一个矩形障碍的房间里规划。论证：若把它当“点 + 把障碍按 Minkowski 差膨胀”，规划结果与原问题等价（提示：纯平移时 $\mathcal{C}_{obs}=\mathcal{O}\ominus\mathcal{A}(\mathbf{0})$）。
\end{enumerate}
\end{practice}

% ════════════════════════════════════════════════════════════════
\section{搜索式规划：图搜索与动态规划}
\label{sec:plan-search}

上一节把规划问题统一成“在 $\mathcal{C}_{free}$ 中连点找路”。最直接的求解思路是\emph{离散化}：把 $\mathcal{C}_{free}$ 划成栅格或图，再在图上搜索。本节先把图搜索的通用框架讲清，下一节专攻 A$^*$ 的最优性。

\paragraph{动机：先把连续问题离散成图。}
把 $\mathcal{C}_{free}$ 按分辨率 $k_i$ 划成栅格，栅格点 $\mathbf{q}=\sum_i j_i\Delta\mathbf{q}_i$（$\Delta\mathbf{q}_i$ 是第 $i$ 轴步长向量），相邻格用边连接（1-邻域 $N_1$ 至多 $2n$ 个邻居、$n$-邻域至多 $3^n-1$ 个），边的碰撞在搜索中按需检验（不预构 $\mathcal{C}_{obs}$）。于是规划变成离散图搜索。下面先在抽象的\emph{离散状态空间}上把搜索讲透。

\begin{definition}[离散可行规划问题]\label{def:plan-discrete}
给定 (1) 非空状态空间 $X$（有限或可数）；(2) 每个 $\mathbf{x}\in X$ 的有限动作空间 $U(\mathbf{x})$；(3) 状态转移函数 $f$，状态转移方程 $\mathbf{x}'=f(\mathbf{x},u)$；(4) 初始状态 $\mathbf{x}_I$；(5) 目标集 $X_G\subset X$。任务：找有限动作序列把 $\mathbf{x}_I$ 变换到某 $X_G$ 中状态。可用\emph{有向状态转移图}表示：顶点 $=X$，有向边 $\mathbf{x}\to\mathbf{x}'$ 当存在 $u$ 使 $\mathbf{x}'=f(\mathbf{x},u)$。
\end{definition}

\noindent 状态转移图常\emph{不显式给出}而在搜索中增量揭示（如魔方：每态 12 动作，目标=六面纯色，规则隐式编码巨大图）。

\subsection{通用前向搜索框架}
\label{sec:plan-forward}

搜索中任一时刻，状态分三类\cite{lavalle2006planning}：\textbf{未访问}（尚未遇到，初始除 $\mathbf{x}_I$ 外全部）；\textbf{死}（已访问且其全部后继都已访问，对搜索无新贡献）；\textbf{活}（已遇到但可能仍有未访问后继，初始唯一活态是 $\mathbf{x}_I$）。活态存于优先队列 $Q$。

\begin{algo}[通用前向搜索]\label{alg:plan-forward}
所有图搜索算法都是下面框架的特例——\textbf{唯一实质的区别，是对 $Q$ 排序的那个函数}。
\end{algo}

\begin{codebox}[text]{FORWARD-SEARCH（LaValle 图 2.4）}
FORWARD_SEARCH
 1  Q.Insert(x_I) and mark x_I as visited
 2  while Q not empty do
 3      x <- Q.GetFirst()
 4      if x in X_G then return SUCCESS
 5      forall u in U(x)
 6          x' <- f(x, u)
 7          if x' not visited
 8              Mark x' as visited
 9              Q.Insert(x')
10          else
11              Resolve duplicate x'          // 见各算法; Dijkstra/A* 在此松弛
12  return FAILURE
\end{codebox}

\noindent 实现要点：第 6 行后记 $\mathbf{x}'$ 的父指针 $\mathbf{x}$，终态回溯到 $\mathbf{x}_I$ 即得路径；判重（第 7 行）树形图无需、栅格用查找表 $O(1)$、一般情形须与 $Q$ 中及死态比较（用哈希缓解）。各具体方法仅 $Q$ 排序不同：

\begin{itemize}
  \item \textbf{广度优先 (BFS)}：$Q$ = FIFO。前沿均匀扩张，所有 $k$ 步计划在 $k+1$ 步之前穷尽，故首个解步数最少。运行 $O(|V|+|E|)$，系统（有限/无穷图均访问每个可达态）。
  \item \textbf{深度优先 (DFS)}：$Q$ = 栈（LIFO）。激进深入，偏好早期长计划（具体哪条任意）。有限 $X$ 系统、无穷 $X$ 不系统。运行 $O(|V|+|E|)$。
  \item \textbf{Dijkstra}：$Q$ 按到达代价排序（\cref{sec:plan-dijkstra}）。
  \item \textbf{A$^*$}：$Q$ 按 $g+h$ 排序（\cref{sec:plan-astar}）。
  \item \textbf{最佳优先 (best-first)}：$Q$ 仅按 cost-to-go 估计排序，贪心、不保最优、不系统（\cref{sec:plan-astar} 反例）。
  \item \textbf{迭代加深 / IDA$^*$}：用 DFS 逐层放宽深度上限（或 $g+h$ 上限），把 DFS 变系统、省空间，分支因子大时优选。
\end{itemize}

\paragraph{后向与双向搜索。}
亦可从 $\mathbf{x}_G$ 反向搜（用反向转移 $\mathbf{x}=f^{-1}(\mathbf{x}',u^{-1})$，等价于在反图上前向搜）；当分支因子在 $\mathbf{x}_I$ 端大时更高效。\textbf{双向搜索}：一树从 $\mathbf{x}_I$、一树从 $\mathbf{x}_G$ 同时生长，相遇即成功，常大幅减少探索面积。这一“双向”思想在采样式 RRT-Connect（\cref{sec:plan-rrtconnect}）里再现并威力更大。

\subsection{动态规划与 Bellman 递推}
\label{sec:plan-dp}

Dijkstra/A$^*$ 的最优性建立在动态规划之上。考虑\emph{最优}规划：每边代价 $l(\mathbf{x},u)\ge0$，最小化阶段可加损失泛函 $L(\pi_K)=\sum_{k=1}^K l(\mathbf{x}_k,u_k)+l_F(\mathbf{x}_F)$（$l_F=0$ 若到达目标、否则 $\infty$，把可行性转成优化）。

\begin{theorem}[最优性原理与 Bellman 递推]\label{thm:plan-bellman}
定义最优 cost-to-go $G_k^*(\mathbf{x}_k)=\min_{u_k,\dots,u_K}\{\sum_{i=k}^K l(\mathbf{x}_i,u_i)+l_F(\mathbf{x}_F)\}$。则它满足后向递推
\begin{equation}
  G_k^*(\mathbf{x}_k)=\min_{u_k}\big\{l(\mathbf{x}_k,u_k)+G_{k+1}^*(\mathbf{x}_{k+1})\big\},\qquad\mathbf{x}_{k+1}=f(\mathbf{x}_k,u_k).
  \label{eq:plan-bellman}
\end{equation}
\end{theorem}
\begin{proof}
把 $G_k^*$ 的定义中第一项 $l(\mathbf{x}_k,u_k)$ 拆出、把对 $u_{k+1},\dots,u_K$ 的内层 $\min$ 分离：
\[
  G_k^*(\mathbf{x}_k)=\min_{u_k}\Big\{l(\mathbf{x}_k,u_k)+\min_{u_{k+1},\dots,u_K}\big[\textstyle\sum_{i=k+1}^K l(\mathbf{x}_i,u_i)+l_F(\mathbf{x}_F)\big]\Big\}.
\]
内层 $\min$ 恰是 $G_{k+1}^*(\mathbf{x}_{k+1})$（最优 cost-to-go 不依赖到达 $\mathbf{x}_{k+1}$ 的历史，此即\emph{最优性原理}），代入即 \cref{eq:plan-bellman}。
\end{proof}

\noindent 当损失 $l\ge0$ 时，无界后向值迭代 $G_F^*\to G_K^*\to\cdots$ 必收敛到\textbf{稳态}，去掉阶段下标得 Bellman 方程 $G^*(\mathbf{x})=\min_u\{l(\mathbf{x},u)+G^*(f(\mathbf{x},u))\}$；从 $G^*$ 逐步取 $u^*=\arg\min_u\{l(\mathbf{x},u)+G^*(f(\mathbf{x},u))\}$ 即得最优动作序列。$G^*$ 是把系统从任意态最优引向目标的\emph{导航函数}。\textbf{负环}会使代价 $\to-\infty$，须禁（可预检）。

\begin{example}[五态后向值迭代]\label{ex:plan-vi}
取 $X=\{a,b,c,d,e\}$，$K=4$，$\mathbf{x}_I=a$，$X_G=\{d\}$。按 \cref{eq:plan-bellman} 从 $G_5^*=l_F$（$d$ 处 0、余 $\infty$）逐阶回推，得 cost-to-go 表（值越小越好）：
\begin{center}\small
\begin{tabular}{l|ccccc}
\toprule
 & $a$ & $b$ & $c$ & $d$ & $e$ \\
\midrule
$G_5^*$ & $\infty$ & $\infty$ & $\infty$ & $0$ & $\infty$ \\
$G_4^*$ & $\infty$ & $4$ & $1$ & $\infty$ & $\infty$ \\
$G_3^*$ & $6$ & $2$ & $\infty$ & $2$ & $\infty$ \\
$G_2^*$ & $4$ & $6$ & $3$ & $\infty$ & $\infty$ \\
$G_1^*$ & $6$ & $4$ & $5$ & $4$ & $\infty$ \\
\bottomrule
\end{tabular}
\end{center}
$G_1^*(a)=6$ 即从 $a$ 出发 4 步内到目标的最优代价。每阶 $O(|X||U|)$，全程 $O(K|X||U|)$。读者可对照验证每个表项都由上一行经 \cref{eq:plan-bellman} 取最小得出。
\end{example}

\subsection{Dijkstra 算法}
\label{sec:plan-dijkstra}

每边代价 $l(e)\ge0$。$Q$ 按\textbf{到达代价}（cost-to-come）$C:X\to[0,\infty]$ 排序；$C^*(\mathbf{x})$ 是从 $\mathbf{x}_I$ 到 $\mathbf{x}$ 的最优到达代价，未知最优时记 $C(\mathbf{x})$。初始 $C^*(\mathbf{x}_I)=0$。每生成 $\mathbf{x}'$ 算 $C(\mathbf{x}')=C^*(\mathbf{x})+l(\mathbf{x},u)$（\cref{alg:plan-forward} 第 11 行：若 $\mathbf{x}'$ 已在 $Q$ 中且新路径更优，则下调 $C(\mathbf{x}')$ 并重排——此即\emph{松弛}）。

\begin{codebox}[text]{Dijkstra（优先队列版，含松弛）}
DIJKSTRA(Graph, source)
 1  dist[source] <- 0; dist[v] <- +inf for v != source
 2  Q <- min-priority queue of all vertices keyed by dist
 3  while Q not empty:
 4      u <- Q.extract_min()
 5      for each neighbor v of u still in Q:
 6          alt <- dist[u] + edge_cost(u, v)          // relaxation
 7          if alt < dist[v]:
 8              dist[v] <- alt; prev[v] <- u
 9              Q.decrease_priority(v, alt)
10  return dist, prev
\end{codebox}

\begin{theorem}[Dijkstra 出队即最优]\label{thm:plan-dijkstra}
当状态 $\mathbf{x}$ 经 \texttt{Q.extract\_min()} 出队（变死态）时，$C(\mathbf{x})=C^*(\mathbf{x})$——即不可能再以更低代价到达 $\mathbf{x}$。
\end{theorem}
\begin{proof}（归纳，LaValle\cite{lavalle2006planning} §2.2.2）
\textbf{基例}：$C^*(\mathbf{x}_I)=0$ 已知。\textbf{归纳假设}：每个死态的最优到达代价已正确确定且不再变。对 $Q$ 的队首 $\mathbf{x}$（其 $C(\mathbf{x})$ 在 $Q$ 中最小），断言 $C(\mathbf{x})=C^*(\mathbf{x})$：任何\emph{总代价更低}的到达路径都必须经过 $Q$ 中另一状态（因起点已死、要走出死态区只能经活态），而那些活态的 $C$ 值\emph{不低于} $\mathbf{x}$（$\mathbf{x}$ 是队首/最小），其后的边代价非负只会让总代价更大；而只经死态的路径已在计算 $C(\mathbf{x})$ 时考虑过。故 $C(\mathbf{x})$ 已最优。探索 $\mathbf{x}$ 全部出边后 $\mathbf{x}$ 变死，归纳继续。
\end{proof}

\begin{pitfall}[概念误区：以为 Dijkstra 适用负权]
\cref{thm:plan-dijkstra} 的证明用到“经过 $Q$ 中其它状态的路径代价不低于队首”，这\emph{依赖边代价非负}。若有负权，一条先经高代价态、再走负边的路径可能反超队首，破坏“出队即最优”的不变量；负环更使代价无下界。负权须改用 Bellman--Ford。运动规划里代价通常是几何长度或能量，天然非负，故 Dijkstra/A$^*$ 适用。
\end{pitfall}

\noindent\textbf{复杂度}：用 Fibonacci 堆实现 $Q$（\texttt{decrease-key} 摊还 $O(1)$）得 $O(|E|+|V|\log|V|)$，是任意非负权有向图单源最短路的渐近最快已知算法；二叉堆实现为 $O((|E|+|V|)\log|V|)$。Dijkstra 属更广的\emph{标签校正}族——允许死态在发现更优代价时复活；它每态只访问一次，故是该族的特例。

\begin{practice}
\begin{enumerate}
  \item 在 \cref{ex:plan-vi} 的图上跑前向 Dijkstra（从 $a$），逐步写出每次出队的状态与其最终 $C^*$，对照后向值迭代的 $G_1^*$。
  \item 构造一个含\emph{负边}（无负环）的小图，说明 Dijkstra 给出错误结果而 \cref{thm:plan-dijkstra} 的归纳在哪一步失效。
\end{enumerate}
\end{practice}

% ════════════════════════════════════════════════════════════════
\section{A$^*$ 搜索与启发函数}
\label{sec:plan-astar}

\paragraph{动机：Dijkstra 太“盲目”。}
Dijkstra 向四面八方均匀扩展（$f=g$），在大地图上会探索大量与目标方向无关的状态。若我们能\emph{预估}“从某态到目标还要多少代价”，就能把搜索“拉向”目标，大幅减少探索——这正是 A$^*$ 的思想。A$^*$ 由 Hart、Nilsson、Raphael（斯坦福研究院）1968 年为移动机器人 Shakey 项目提出\cite{hart1968astar}；和式 $g+h$ 由 Raphael 建议，可采纳性与一致性概念由 Hart 引入。

\begin{definition}[A$^*$ 评价函数]\label{def:plan-astar}
每节点 $\mathbf{x}$ 维护
\begin{equation}
  f(\mathbf{x})=g(\mathbf{x})+h(\mathbf{x}),
  \label{eq:plan-fgh}
\end{equation}
其中 $g(\mathbf{x})$ = 起点到 $\mathbf{x}$ 的当前已知最小代价（cost-to-come），$h(\mathbf{x})$ = $\mathbf{x}$ 到目标的启发式估计代价（cost-to-go 的下估计），$f(\mathbf{x})$ = “若路径经 $\mathbf{x}$，从起到终最廉价代价的当前最佳猜测”。\textbf{A$^*$ 与 Dijkstra 的唯一区别就是把 $Q$ 排序键从 $g$ 换成 $f=g+h$}；当 $h\equiv0$ 时 A$^*$ 退化为 Dijkstra。
\end{definition}

\begin{codebox}[text]{A*（Wikipedia/HNR68 结构）}
A_STAR(start, goal, h)
 1  openSet <- {start}
 2  gScore[start] <- 0; gScore[n] <- +inf for n != start
 3  fScore[start] <- h(start)
 4  while openSet not empty:
 5      current <- node in openSet with lowest fScore       // (1) f-min
 6      if current = goal: return reconstruct_path(current) // (2)
 7      openSet.remove(current)
 8      for each neighbor of current:
 9          tentative_g <- gScore[current] + d(current, neighbor)
10          if tentative_g < gScore[neighbor]:              // (3) relax
11              cameFrom[neighbor] <- current
12              gScore[neighbor]   <- tentative_g
13              fScore[neighbor]   <- tentative_g + h(neighbor)
14              if neighbor not in openSet: openSet.add(neighbor)
15  return FAILURE
\end{codebox}

\subsection{可采纳性与最优性}
\label{sec:plan-admissible}

\begin{definition}[可采纳启发]\label{def:plan-admissible}
启发 $h$ \textbf{可采纳}（admissible），若对所有 $\mathbf{x}$ 都\emph{从不高估}到目标的真实最优代价 $h^*(\mathbf{x})$：
\begin{equation}
  h(\mathbf{x})\le h^*(\mathbf{x})\quad\forall\mathbf{x}.
  \label{eq:plan-admissible}
\end{equation}
\end{definition}

\noindent 构造可采纳启发的标准手法是\textbf{松弛问题}：去掉部分约束后的精确解代价 $\le$ 原问题，故为下估计。例（栅格迷宫，代价=步数）：曼哈顿距离 $|i'-i|+|j'-j|$ 是“忽略障碍的直达计划长度”，加入障碍后绕行只会使代价增加，故可采纳；$h\equiv0$ 也可采纳但无信息（退化为 Dijkstra）。又如 8-数码：错位棋子数（每个不在位的至少移一次）、各棋子曼哈顿距离之和，都可采纳。

\begin{theorem}[A$^*$ 在可采纳启发下最优]\label{thm:plan-astar-opt}
若 $h$ 可采纳，则 A$^*$（树搜索）返回的解是最优解。
\end{theorem}
\begin{proof}（反证，Russell--Norvig 标准证明\rebuilt——据 Wikipedia 直觉链 + AIMA §3.5.2 补全字句）
设最优解代价为 $C^*$。反设 A$^*$ 选中一个\emph{次优}目标 $G_2$，即 $g(G_2)>C^*$（目标处 $h(G_2)=0$ 故 $f(G_2)=g(G_2)>C^*$）。在 A$^*$ 即将选中 $G_2$ 的那一刻，最优路径上必有某节点 $\mathbf{n}$ 仍在 open 中（起点在最优路径上、$G_2$ 不在最优路径上，故最优路径未被完全展开）。对该 $\mathbf{n}$：
\[
  f(\mathbf{n})=g(\mathbf{n})+h(\mathbf{n})\le g(\mathbf{n})+h^*(\mathbf{n})=C^*,
\]
第一不等号由可采纳 \cref{eq:plan-admissible}，末等号因 $g(\mathbf{n})+h^*(\mathbf{n})$ 恰是过 $\mathbf{n}$ 的最优路径代价 $=C^*$。于是 $f(\mathbf{n})\le C^*<f(G_2)$。但 A$^*$ 总取 open 中 $f$ \emph{最小}者，应先扩展 $\mathbf{n}$ 而非 $G_2$——矛盾。故 A$^*$ 不会返回次优解。
\end{proof}

\begin{pitfall}[思维陷阱：best-first 看着像 A$^*$ 却不保最优]
最佳优先搜索（greedy best-first）只按 cost-to-go 估计 $h$ 排序（不计已走代价 $g$），故\emph{不在意}是否高估，解不必最优。它常太贪心、被早期“看起来好”的态误导（LaValle\cite{lavalle2006planning} 的 3D 螺旋管反例：从管口进入会绕整条螺旋而非直奔目标，浪费任意多步）。best-first 不系统，最坏比 A$^*$/Dijkstra 更差。A$^*$ 与 best-first 的本质区别就在于 A$^*$ \emph{计入了 $g(\mathbf{x})$}。
\end{pitfall}

\subsection{一致启发与无重复展开}
\label{sec:plan-consistent}

可采纳保证树搜索最优，但在\emph{图搜索}（关闭重复状态）中，仅可采纳\emph{不足}——若 $h$ 可采纳却不一致，一个节点每次以更低的“目前最佳”代价被发现时需重新展开（reopen），既慢又易错。要在图搜索中无重复地保最优，需更强的\emph{一致性}。

\begin{definition}[一致（单调）启发]\label{def:plan-consistent}
启发 $h$ \textbf{一致}（consistent / monotone），若对每条边 $(\mathbf{x},\mathbf{y})$ 满足“三角不等式”，且目标处为零：
\begin{equation}
  h(\mathbf{x})\le c(\mathbf{x},\mathbf{y})+h(\mathbf{y}),\qquad h(G)=0.
  \label{eq:plan-consistent}
\end{equation}
\end{definition}

\begin{theorem}[一致蕴含可采纳]\label{thm:plan-cons-adm}
一致启发必可采纳。
\end{theorem}
\begin{proof}（归纳）
设目标 $\mathbf{N}_0$ 处 $h(\mathbf{N}_0)=0$，是基例（$0\le0$）。沿任一到目标的路径 $\mathbf{N}_{i+1}\to\mathbf{N}_i\to\cdots\to\mathbf{N}_0$ 反复用一致性 \cref{eq:plan-consistent}：
\[
  h(\mathbf{N}_{i+1})\le c(\mathbf{N}_{i+1},\mathbf{N}_i)+h(\mathbf{N}_i)\le\cdots\le\sum_{j} c(\mathbf{N}_{j+1},\mathbf{N}_j)+h(\mathbf{N}_0).
\]
右端 $=$ 从 $\mathbf{N}_{i+1}$ 沿该路径到目标的实际代价（$h(\mathbf{N}_0)=0$）。对\emph{任意}到目标的路径都成立，取最优路径即得 $h(\mathbf{N}_{i+1})\le h^*(\mathbf{N}_{i+1})$，即可采纳。
\end{proof}

\begin{theorem}[一致蕴含 $f$ 沿路径单调非减、出队即最优]\label{thm:plan-cons-mono}
一致启发下，$f$ 沿任一路径单调非减；故 A$^*$ 弹出（扩展）一个节点时其 $g$ 值已最优，每节点至多处理一次。
\end{theorem}
\begin{proof}
设 $\mathbf{y}$ 为 $\mathbf{x}$ 的后继。则
\[
  f(\mathbf{y})=g(\mathbf{y})+h(\mathbf{y})=g(\mathbf{x})+c(\mathbf{x},\mathbf{y})+h(\mathbf{y})\ \ge\ g(\mathbf{x})+h(\mathbf{x})=f(\mathbf{x}),
\]
不等号正是一致性 $c(\mathbf{x},\mathbf{y})+h(\mathbf{y})\ge h(\mathbf{x})$。故 $f$ 单调非减。又由单调性，A$^*$ 按 $f$ 升序弹出节点，弹出时到该节点的最优路径已被找到（其 $g$ 已最优，同 \cref{thm:plan-dijkstra} 的论证），故无需 reopen。
\end{proof}

\begin{insight}
一致启发下 A$^*$ \textbf{等价于在“约化代价” $c'(\mathbf{x},\mathbf{y})=c(\mathbf{x},\mathbf{y})-h(\mathbf{x})+h(\mathbf{y})\ge0$ 上跑 Dijkstra}（约化代价非负正是一致性）。这把 A$^*$ 严格归约为 Dijkstra，也解释了为何 Dijkstra（$h\equiv0$，平凡一致）是 A$^*$ 的特例。所有一致启发都可采纳，但反之不真——实践中由松弛问题/三角不等式得到的启发（曼哈顿、欧氏、对角）通常既可采纳又一致。\textbf{最优效率}（Dechter--Pearl）：给定一致启发，A$^*$ 在所有可采纳的同类算法中扩展节点数最少（除平局处理外）。
\end{insight}

\subsection{栅格地图的启发函数}
\label{sec:plan-grid}

栅格地图把 $\mathcal{C}_{free}$ 离散为方格，相邻格连边后跑 A$^*$。设 $D$ = 正交相邻格移动代价，$\mathrm{dx}=|x_G-x_{\mathbf{x}}|$，$\mathrm{dy}=|y_G-y_{\mathbf{x}}|$。常用启发（Amit Patel/斯坦福\cite{patel2018heuristics}）：

\begin{align}
  \text{曼哈顿（4-邻接）：}&\quad h(\mathbf{x})=D\,(\mathrm{dx}+\mathrm{dy}); \label{eq:plan-manhattan}\\
  \text{对角（8-邻接，斜走代价 }D_2\text{）：}&\quad h(\mathbf{x})=D\,\max(\mathrm{dx},\mathrm{dy})+(D_2-D)\,\min(\mathrm{dx},\mathrm{dy}); \label{eq:plan-diagonal}\\
  \text{欧氏（任意角度）：}&\quad h(\mathbf{x})=D\sqrt{\mathrm{dx}^2+\mathrm{dy}^2}. \label{eq:plan-euclid}
\end{align}
\cref{eq:plan-diagonal} 的特例：Chebyshev（$D=1,D_2=1$）得 $h=\max(\mathrm{dx},\mathrm{dy})$；octile（$D=1,D_2=\sqrt2$）得 $h=\max(\mathrm{dx},\mathrm{dy})+(\sqrt2-1)\min(\mathrm{dx},\mathrm{dy})$。

\begin{pitfall}[编程/概念陷阱：用平方欧氏距离当启发]
\textbf{错误做法}：写 $h=D(\mathrm{dx}^2+\mathrm{dy}^2)$（省一个开方）。\textbf{现象}：A$^*$ 路径变形、慢甚至非最优。\textbf{根本原因}：当 A$^*$ 算 $f=g+h$ 时，距离的\emph{平方}量纲与代价 $g$（线性距离）不匹配，远处节点的 $h$ 被严重\emph{高估}，破坏可采纳性。\textbf{正确做法}：启发返回的单位（米、步、秒）必须与 $g$ 一致；要省开方就用曼哈顿/octile（线性、可采纳）。\textbf{自检}：检查 $h$ 在“直线无障碍”情形是否恰等于真代价、是否处处 $\le h^*$。
\end{pitfall}

\begin{pitfall}[思维陷阱：欠估/高估的权衡与平局摊开]
$h$ 欠估（含 $h\equiv0$=Dijkstra）保证最优但扩展更多节点（更慢）；$h$ 恰等 $h^*$ 则 A$^*$ 只沿最优路径扩展（最快，但一般不可得）；$h$ 高估则放弃最优换速度。另一个隐蔽问题：当多节点 $f$ 相等时 A$^*$ 路径会“摊开”成方块状（无谓探索）。对策：对启发乘一极小因子偏向目标 $h\mathrel{*}=(1+p)$，$p\approx1/(\text{预期最大路径长})$，$p$ 足够小则仍近似可采纳，却能显著减少等价 $f$ 的探索。
\end{pitfall}

\paragraph{$\dagger$ D$^*$：动态/部分已知环境的增量 A$^*$。}\rebuilt
A$^*$ 假设地图完整且静止。Stentz 1994 的 D$^*$（Dynamic A$^*$）允许弧代价在运行中改变\cite{stentz1994optimal}：移动机器人先按已知（或乐观假设）地图规划，行进中传感器发现新障碍/代价变化，需\emph{高效}修正而非整路重算。D$^*$ \textbf{反向搜索}（树根在 goal），对每个可能的 start 都算出到 goal 的最优路径，故机器人移动后无需重启。其关键是一个\emph{历史最小} key
\begin{equation}
  k(X)=\min\big(h(X),\ X\text{ 入 OPEN 以来 }h(X)\text{ 的历史值}\big),
  \label{eq:plan-dstar-key}
\end{equation}
（这里 $h(X)$ 是“到目标代价”cost-to-goal，\emph{不是} A$^*$ 的启发下界——记号撞名须当心）。key 把 OPEN 状态分两类：\textbf{LOWER} 状态（$k=h$，代价下降，向邻居传播更低代价）与 \textbf{RAISE} 状态（$k<h$，代价曾被抬高，向邻居传播上升并寻找更优新父）。核心函数 \texttt{PROCESS-STATE} 弹出 key 最小态、按 RAISE/LOWER 双波在图上传播代价变化；\texttt{MODIFY-COST} 在检测到弧代价变化时把受影响态放回 OPEN 触发局部重传播。其关键性质（逻辑等价复建，原文为扫描件无法逐字复现）：当状态被置 CLOSED 且 $k_{\min}\ge h(X)$ 时 $h(X)$ 即当前已知代价下到 goal 的\emph{最优}值（同 Dijkstra 出队即最优）；弧代价改变后 D$^*$ 的结果与“用新代价从头重算”\emph{相同}，但只触及受影响子集，故远快于整图重算。现代实践首选其后继 \textbf{D$^*$ Lite}（Koenig \& Likhachev 2002，基于 LPA$^*$ 重构，逻辑更简、性能相当或更优）。

\begin{practice}
\begin{enumerate}
  \item 证明曼哈顿距离 \cref{eq:plan-manhattan}（$D=1$，4-邻接栅格）一致：验证对每条单位边 $h(\mathbf{x})\le1+h(\mathbf{y})$。
  \item （开放思考）若把对角启发 \cref{eq:plan-diagonal} 用在只许 4-邻接移动的栅格上，它还可采纳吗？给出论证或反例（提示：比较斜向真代价与 $D_2$）。
  \item 在含一个新出现障碍的栅格上手工跑一遍 D$^*$ 的 RAISE/LOWER 波：标出哪些格的 $h$ 被抬高、代价波如何局部传播，对照“整图重跑 A$^*$”确认结果一致。
\end{enumerate}
\end{practice}

% ════════════════════════════════════════════════════════════════
\section{采样式运动规划：RRT 与 PRM}
\label{sec:plan-sampling}

\paragraph{动机：栅格搜索撞上维数灾难。}
上两节的栅格 + A$^*$ 在低维（平面导航）极有效，但 \cref{thm:plan-pspace} 与栅格法 $O(1/h^d)$ 的复杂度（$h$ 分辨率、$d$ 维数）说明：六自由度机械臂若每轴取 100 个格点，就有 $100^6=10^{12}$ 个格点，遑论显式建 $\mathcal{C}_{obs}$。\textbf{反面}：高维栅格搜索还易困在局部极小（如 $\mathcal{C}_{obs}$ 的“碗”会先被 $100^n$ 个格点填满）。采样式规划绕开这两难：\emph{回避显式 $\mathcal{C}_{obs}$}，用随机采样探测 $\mathcal{C}_{free}$、把碰撞检测当黑箱，再用前几节的图搜索取路径。

\paragraph{完备性的弱化。}
精确几何方法\textbf{完备}（有限时间正确报告有无解，\cref{thm:plan-pspace} 的代价使其不实用）。采样式只能达到弱化概念：\textbf{分辨率完备}（确定性稠密采样——分辨率够细则有解必找到，无解可能永跑）；\textbf{概率完备}（随机采样——样本足够多时找到存在解的概率趋 1，见 \cref{def:plan-probcomplete}）。

\subsection{通用原语}
\label{sec:plan-primitives}

所有采样算法共用以下原语\cite{karaman2011sampling}（构型空间取 $X=(0,1)^d$，$d\ge2$）：
\begin{itemize}
  \item \textbf{SampleFree}：返回 $\mathcal{C}_{free}$ 上独立同分布（默认均匀）样本序列 $\{\mathbf{x}_i\}$。
  \item \textbf{Nearest}$(G,\mathbf{x})=\arg\min_{\mathbf{v}\in V}\|\mathbf{x}-\mathbf{v}\|$（欧氏最近顶点）；集值版 \textbf{kNearest}$(G,\mathbf{x},k)$ 取最近 $k$ 个。
  \item \textbf{Near}$(G,\mathbf{x},r)=\{\mathbf{v}\in V:\mathbf{v}\in\mathcal{B}_{\mathbf{x},r}\}$（半径 $r$ 球内顶点）。
  \item \textbf{Steer}$(\mathbf{x},\mathbf{y})=\arg\min_{\mathbf{z}\in\mathcal{B}_{\mathbf{x},\eta}}\|\mathbf{z}-\mathbf{y}\|$：返回从 $\mathbf{x}$ 朝 $\mathbf{y}$ 前进至多 $\eta$ 的点（$\eta>0$ 预设步长）。
  \item \textbf{CollisionFree}$(\mathbf{x},\mathbf{x}')$：线段 $[\mathbf{x},\mathbf{x}']\subset\mathcal{C}_{free}$ 则真。
\end{itemize}
所有算法输出图 $G=(V,E)$（$V\subset\mathcal{C}_{free}$，$|V|\le n+1$），解由标准最短路（Dijkstra）从图中提取——这就是为什么采样式被称为“随机离散化 + 图搜索”。

\subsection{RRT：快速探索随机树}
\label{sec:plan-rrt}

RRT（Rapidly-exploring Random Tree，LaValle 1998 / Kuffner--LaValle 2000\cite{kuffner2000rrtconnect}）增量地建一棵以 $\mathbf{q}_I$ 为根的可行轨迹树，\emph{无需参数调校}即有好表现。每轮采一个随机样本，把树中最近顶点朝它扩一步。

\begin{algo}[RRT（含障碍）]\label{alg:plan-rrt}
\end{algo}
\begin{codebox}[text]{RRT（Karaman-Frazzoli Alg.3 / Kuffner-LaValle EXTEND）}
RRT
 1  V <- {x_init};  E <- {};
 2  for i = 1, ..., n do
 3      x_rand    <- SampleFree_i;
 4      x_nearest <- Nearest(G=(V,E), x_rand);
 5      x_new     <- Steer(x_nearest, x_rand);          // advance at most eta
 6      if CollisionFree(x_nearest, x_new) then
 7          V <- V union {x_new};  E <- E union {(x_nearest, x_new)};
 8  return G = (V, E);
\end{codebox}

\noindent 原始 RRT 的 \texttt{EXTEND} 有三种返回：\textbf{Reached}（$\mathbf{q}_{rand}$ 在 $\eta$ 内，直接加入）、\textbf{Advanced}（新顶点 $\mathbf{q}_{new}\ne\mathbf{q}_{rand}$ 加入）、\textbf{Trapped}（提出的新顶点不在 $\mathcal{C}_{free}$，拒绝）。求解查询时常用\textbf{目标偏置}：以小概率 $p$（经验 $1/100$）直接把 $\mathbf{q}_{rand}$ 设为目标，引导树朝目标生长（偏置太强像贪心势场、太弱无动力连目标）。

\begin{insight}
\textbf{RRT 为何“快速探索”？——Voronoi 偏置。}\;
\texttt{Nearest} 把新样本连到树上最近点，\emph{等价于}“某顶点被选中扩展的概率正比于其 Voronoi 区域的体积”。大 Voronoi 区域 = 大片未探索空间，故树被“拉”向边界/未探索区，呈分形式快速铺满（LaValle 图 5.19：45 次迭代即触达四角，2345 次渐稠密）。极限下树以概率 1 \emph{稠密}（任意构型可被任意接近）——这是 RRT 概率完备的几何根。
\end{insight}

\begin{figure}[ht]
\centering
\small
\begin{tikzpicture}[>=Stealth, scale=1.0]
  % bounding box
  \draw[gray!50] (0,0) rectangle (5,3.4);
  % root
  \coordinate (root) at (1,1);
  \fill[main!70!black] (root) circle (1.6pt) node[below left, font=\scriptsize]{$\mathbf{q}_I$};
  % a sparse fractal-like tree (illustrative)
  \coordinate (a) at (1.9,1.6); \coordinate (b) at (2.7,0.7);
  \coordinate (c) at (3.0,2.4); \coordinate (d) at (4.1,2.9);
  \coordinate (e) at (4.3,1.1); \coordinate (g) at (2.1,2.8);
  \coordinate (f) at (0.4,2.5);
  \foreach \p/\q in {root/a, root/b, root/f, a/c, a/g, c/d, b/e}
    \draw[main!60!black, thick] (\p)--(\q);
  \foreach \p in {a,b,c,d,e,f,g} \fill[main!60!black] (\p) circle (1.2pt);
  % a random sample and the extension toward it
  \coordinate (rand) at (4.6,3.1);
  \fill[second!70!black] (rand) circle (1.4pt) node[above, font=\scriptsize]{$\mathbf{q}_{rand}$};
  \draw[second!70!black, dashed, ->] (d) -- (rand) node[midway, right=1pt, font=\scriptsize]{Steer $\eta$};
  \node[font=\scriptsize, gray] at (2.5,0.2) {树偏向大 Voronoi 区（未探索处）扩展};
\end{tikzpicture}
\caption{RRT 的快速探索与 Voronoi 偏置：新样本 $\mathbf{q}_{rand}$ 连到树中最近点并经 \texttt{Steer} 前进 $\eta$；选中扩展的概率正比于顶点 Voronoi 区面积，故树优先铺向大片未探索空间（仿 LaValle\cite{lavalle2006planning} 图 5.19 风格重绘）。}
\label{fig:plan-rrt}
\end{figure}

\paragraph{RDT 视角与无障碍变体。}
LaValle 把 RRT 归入\emph{快速探索稠密树}（RDT）族：序列随机即 RRT、确定性即 RDT。无障碍时直接连到“树已达点集”swath $S=\bigcup_{e\in E}e([0,1])$ 中最近点 $\mathbf{q}_n$（若 $\mathbf{q}_n$ 落在某边内部则\emph{裂边}成新顶点再连）；含障碍时用 \texttt{stopping-configuration} 沿 $\mathbf{q}_n\to\mathbf{q}_{rand}$ 方向返回碰到边界前最后可达构型。高效最近邻查找用 Kd-树（$k$ 点建树 $O(nk\lg k)$、查最近 $O(\lg k)$，实用约到 20 维）。

\subsection{RRT-Connect：双树 + 贪心连接}
\label{sec:plan-rrtconnect}

\paragraph{动机：单树连目标慢，双树相向更快。}
单棵 RRT 要靠目标偏置“碰运气”连到目标，在狭窄通道（bug trap）里很慢。Kuffner--LaValle 2000\cite{kuffner2000rrtconnect} 的 RRT-Connect 用两棵树（一从 $\mathbf{q}_I$、一从 $\mathbf{q}_G$）相向生长，并用一个\emph{贪心} \texttt{CONNECT} 启发把一棵树尽力连向另一棵的新顶点。

\begin{codebox}[text]{RRT-Connect（Kuffner-LaValle 图 2/5）}
EXTEND(T, q)
 1  q_near <- NEAREST_NEIGHBOR(q, T);
 2  if NEW_CONFIG(q, q_near, q_new) then    // advance step eps, collision-check
 3      T.add_vertex(q_new); T.add_edge(q_near, q_new);
 4      if q_new = q then return Reached;
 5      else return Advanced;
 6  return Trapped;

CONNECT(T, q)
 1  repeat S <- EXTEND(T, q) until S != Advanced;   // greedily extend until Reached/Trapped
 2  return S;

RRT_CONNECT_PLANNER(q_init, q_goal)
 1  T_a.init(q_init);  T_b.init(q_goal);
 2  for k = 1 to K do
 3      q_rand <- RANDOM_CONFIG();
 4      if EXTEND(T_a, q_rand) != Trapped then        // grow T_a one step
 5          if CONNECT(T_b, q_new) = Reached then       // T_b greedily reaches T_a's new node
 6              return PATH(T_a, T_b);                   // two trees meet
 7      SWAP(T_a, T_b);                                  // alternate roles
 8  return Failure;
\end{codebox}

\noindent\textbf{CONNECT 的精髓}：不是只扩一个 $\epsilon$ 步，而是迭代 \texttt{EXTEND} 直到到达目标点或撞障碍。它像人工势场的吸引函数使快速收敛，但结合了 RRT 的均匀探索，避开势场法的局部极小（“吸引盆随树增长而移动”，而势场法吸引盆固定在目标）。关键优势：一次 CONNECT 只调一次 \texttt{NEAREST\_NEIGHBOR} 即可建一条长路径（每个新顶点成为下一个的最近邻）。把 \texttt{EXTEND} 与 \texttt{CONNECT} 互换可调贪心强度——这是 OMPL/MoveIt 默认规划器之一。

\subsection{PRM：概率路标图（多查询）}
\label{sec:plan-prm}

\paragraph{动机：同一环境多次查询，值得预处理。}
RRT 是\emph{单查询}（给一对 $(\mathbf{q}_I,\mathbf{q}_G)$ 直接算）。若机器人与障碍固定、要回答多个查询，则值得大力\emph{预处理}建一张\textbf{路标图}（roadmap），之后每个查询只需把起终点接入图再做图搜索。这是 Kavraki 等的概率路标（PRM）。

\begin{codebox}[text]{PRM 预处理（Karaman-Frazzoli Alg.1 / LaValle BUILD_ROADMAP）}
PRM (preprocessing)
 1  V <- {};  E <- {};
 2  for i = 0, ..., n do
 3      x_rand <- SampleFree_i;
 4      U <- Near(G=(V,E), x_rand, r);
 5      V <- V union {x_rand};
 6      foreach u in U, in order of increasing ||u - x_rand||, do
 7          if x_rand and u are NOT in the same connected component of G then
 8              if CollisionFree(x_rand, u) then E <- E union {(x_rand,u),(u,x_rand)};
 9  return G = (V, E);
\end{codebox}

\noindent\textbf{查询阶段}：把 $\mathbf{q}_I,\mathbf{q}_G$ 各连入 $G$，再用\emph{第 \ref{sec:plan-dijkstra} 节的 Dijkstra} 在 $G$ 上取最短路——这正是“构型空间 $\to$ 采样离散化 $\to$ 图搜索”主线的闭环。第 7 行 \texttt{not same\_component}（用 union-find 近 $O(1)$ 判）是为效率：只在不同连通分量间连边，使 PRM 路标成\emph{森林}。邻域 \texttt{Near} 常取最近 $K$ 个（典型 $K=15$）或半径 $r$ 球内全部。

\paragraph{sPRM：可分析的简化版。}
文献分析用的是 \textbf{sPRM}（Kavraki 等 1998）：用 $\mathbf{x}_{init}$ 初始化 $V$、采 $n$ 点、连接距离 $<r$ 的点，但\emph{允许}同分量内连接。无障碍时（$\mathcal{C}_{free}=\mathcal{C}$）sPRM 路标恰是\textbf{随机 $r$-disc 图}——这是与随机几何图理论的关键连接（\cref{sec:plan-optimal} 据此给出临界半径）。实用变体：k-Nearest sPRM（连最近 $k$，无障碍时为随机 $k$-NN 图）；bounded-degree sPRM（$U\leftarrow\texttt{Near}\cap\texttt{kNearest}$ 限度数，典型 $k=20$）；variable-radius sPRM（$r$ 随 $n$ 变——但经典文献对 $r$-$n$ 关系无明确指引，这正是 PRM$^*$ 要补的）。

\begin{pitfall}[思维陷阱：以为“连最近邻”就一定连通]
RRT 可看作“1-最近邻 sPRM”的增量版（也连最近邻），但有一个致命差别：RRT \emph{强制增量连通}（每个新点都接到已有树上），而纯 1-最近邻 sPRM（$k=1$，每点只连其最近邻）\textbf{不概率完备}——含起点的连通分量长度 $\sim n^{-1/d}\to0$ 缩成一点，够不到任何固定距离外的目标，失败概率趋 1（\cref{sec:plan-optimal} 定理）。这说明“连最近邻”这个朴素直觉，离“保证连通到目标”差得很远；连接结构（连几个、是否强制连通）才是完备性的关键。
\end{pitfall}

\begin{practice}
\begin{enumerate}
  \item 在一个有狭窄通道的 2D 环境里，分别画出单树 RRT 与 RRT-Connect 的典型生长过程，说明双树为何更易穿过通道（提示：两波前从两侧向通道相向逼近）。
  \item 解释 PRM 第 7 行 \texttt{not same\_component} 为何用 union-find：若去掉它（即 sPRM）会多做哪些碰撞检测？对解质量有何影响（前瞻 \cref{sec:plan-optimal}：sPRM 渐近最优而 PRM 不）。
  \item 给定均匀随机 $n$ 点，论证其弥散（最大空球半径）量级为 $O((\log n/n)^{1/d})$——这将解释 PRM$^*$ 的连接半径为何取此衰减率。
\end{enumerate}
\end{practice}

% ════════════════════════════════════════════════════════════════
\section{渐近最优采样规划：RRT$^*$、RRG、PRM$^*$}
\label{sec:plan-optimal}

\paragraph{动机：RRT/PRM 能找到路，但不是最优路。}
标准 RRT 与 PRM 概率完备，却\emph{不渐近最优}——它们返回的解代价会收敛到一个\emph{次优}值，不会随采样增多而趋向真正最优 $c^*$。Karaman--Frazzoli 2011\cite{karaman2011sampling}（IJRR，本节定理编号均按此源）提出 PRM$^*$、RRG、RRT$^*$，证明它们\emph{渐近最优}（解代价以概率 1 收敛到 $c^*$），且与原算法\emph{同阶}复杂度。关键洞察是：连接半径必须随样本数按 $(\log n/n)^{1/d}$ 缩放。

\subsection{RRT$^*$：选最优父 + 重接线}
\label{sec:plan-rrtstar}

RRT$^*$ = RRG 去环：保留树结构，但对 $\mathbf{x}_{new}$ 的连接做两步优化，确保每顶点经\emph{最小代价路径}到达。先定义 \texttt{Cost}$(\mathbf{v})$ = 根到 $\mathbf{v}$ 的唯一树路径代价（可加代价时 $\texttt{Cost}(\mathbf{v})=\texttt{Cost}(\texttt{Parent}(\mathbf{v}))+c(\texttt{Line}(\texttt{Parent}(\mathbf{v}),\mathbf{v}))$，根为 0）。

\begin{algo}[RRT$^*$]\label{alg:plan-rrtstar}
\end{algo}
\begin{codebox}[text]{RRT*（Karaman-Frazzoli Alg.6）}
RRT*
 1  V <- {x_init};  E <- {};
 2  for i = 1, ..., n do
 3      x_rand    <- SampleFree_i;
 4      x_nearest <- Nearest(G=(V,E), x_rand);
 5      x_new     <- Steer(x_nearest, x_rand);
 6      if CollisionFree(x_nearest, x_new) then
 7          X_near <- Near(G=(V,E), x_new, min{gamma_RRTstar (log(|V|)/|V|)^(1/d), eta});
 8          V <- V union {x_new};
 9          x_min <- x_nearest;  c_min <- Cost(x_nearest) + c(Line(x_nearest, x_new));
10          foreach x_near in X_near do                 // -- (i) choose minimum-cost parent --
11              if CollisionFree(x_near, x_new) and
12                 Cost(x_near) + c(Line(x_near, x_new)) < c_min then
13                  x_min <- x_near;  c_min <- Cost(x_near) + c(Line(x_near, x_new));
14          E <- E union {(x_min, x_new)};
15          foreach x_near in X_near do                 // -- (ii) rewire the tree --
16              if CollisionFree(x_new, x_near) and
17                 Cost(x_new) + c(Line(x_new, x_near)) < Cost(x_near) then
18                  x_parent <- Parent(x_near);
19                  E <- (E \ {(x_parent, x_near)}) union {(x_new, x_near)};
20  return G = (V, E);
\end{codebox}

\noindent\textbf{两步精髓}：(i) \emph{选最优父}（行 10--14）——从 $X_{near}$ 中选能以最小代价 $c_{min}$ 连到 $\mathbf{x}_{new}$ 的顶点作父；(ii) \emph{重接线}（行 15--19）——若经 $\mathbf{x}_{new}$ 到某 $\mathbf{x}_{near}$ 比其当前路径更省，则改 $\mathbf{x}_{near}$ 的父为 $\mathbf{x}_{new}$（删旧父边、加新边，维持树结构）。

\begin{figure}[ht]
\centering
\small
\begin{tikzpicture}[>=Stealth, scale=1.0, font=\scriptsize]
  % left: choose parent
  \begin{scope}
    \node at (1.4,2.6) {\textbf{(i) 选最优父}};
    \coordinate (r) at (0,0); \fill (r) circle (1.3pt) node[left]{root};
    \coordinate (p1) at (1.0,0.8); \coordinate (p2) at (1.2,-0.4); \coordinate (p3) at (0.5,1.4);
    \foreach \p in {p1,p2,p3} \fill (\p) circle (1.1pt);
    \draw[gray!60] (r)--(p1); \draw[gray!60] (r)--(p3); \draw[gray!60] (p1)--(p2);
    \coordinate (xn) at (2.3,0.5); \fill[second!70!black] (xn) circle (1.4pt) node[right]{$\mathbf{x}_{new}$};
    \draw[dashed, gray] (xn) circle (1.25);
    \draw[main!70!black, very thick, ->] (p2)--(xn);   % chosen min-cost parent
    \draw[second!50, dotted] (p1)--(xn); \draw[second!50, dotted] (p3)--(xn);
    \node[main!60!black] at (2.0,-0.2) {min-cost};
  \end{scope}
  % right: rewire
  \begin{scope}[xshift=5.6cm]
    \node at (1.4,2.6) {\textbf{(ii) 重接线}};
    \coordinate (r) at (0,0); \fill (r) circle (1.3pt) node[left]{root};
    \coordinate (q1) at (1.0,0.8); \coordinate (q2) at (1.9,1.3);
    \coordinate (xn) at (1.6,0.0); \fill[second!70!black] (xn) circle (1.4pt) node[below]{$\mathbf{x}_{new}$};
    \fill (q1) circle (1.1pt); \fill (q2) circle (1.1pt) node[right]{$\mathbf{x}_{near}$};
    \draw[gray!60] (r)--(q1);
    \draw[red!60!black, dashed] (q1)--(q2) node[midway, above, red!50!black]{旧父(贵)};
    \draw[main!70!black, very thick, ->] (xn)--(q2) node[midway, below right, main!60!black]{新父(更省)};
    \draw[gray!60] (r)--(xn);
  \end{scope}
\end{tikzpicture}
\caption{RRT$^*$ 的两步：(i) 从半径 $r$ 邻域 $X_{near}$ 中挑能以最小代价到达 $\mathbf{x}_{new}$ 的顶点作父；(ii) 检查邻域内各 $\mathbf{x}_{near}$ 是否经 $\mathbf{x}_{new}$ 到达更省，若是则改其父为 $\mathbf{x}_{new}$（删旧边、加新边）。重接线打破 RRT“一旦连上永不重连”的冻结，是渐近最优的来由。}
\label{fig:plan-rrtstar}
\end{figure}

\paragraph{RRG 与 PRM$^*$。}
\textbf{RRG}（Rapidly-exploring Random Graph）在 RRT 基础上，每加新点就向半径球内\emph{所有}无碰邻居连边（成图、可含环）；同采样序列下 RRT 树是 RRG 图的子图。\textbf{PRM$^*$} 与 sPRM 唯一区别是把固定半径换成 $r(n)$。三者算法结构对照见 \cref{tab:plan-algos}，半径公式见 \cref{sec:plan-radius}。RRT$^*$ 相比 RRG 省内存（树结构）、易扩展到微分约束、且是 \textbf{anytime} 算法（随时给可行解并渐近收敛到最优）。

\subsection{临界连接半径与其几何来由}
\label{sec:plan-radius}

\begin{definition}[渐近最优的连接半径]\label{def:plan-radius}
PRM$^*$/RRG 用 $r$-disc 半径
\begin{equation}
  r(n)=\gamma_{\mathrm{PRM}}\Big(\frac{\log n}{n}\Big)^{1/d},\qquad
  \gamma_{\mathrm{PRM}}>\gamma_{\mathrm{PRM}}^*=2\Big(1+\tfrac1d\Big)^{1/d}\Big(\frac{\mu(\mathcal{C}_{free})}{\zeta_d}\Big)^{1/d},
  \label{eq:plan-gamma-prm}
\end{equation}
其中 $d=\dim\mathcal{C}$，$\mu(\mathcal{C}_{free})$ = 自由空间体积，$\zeta_d=\pi^{d/2}/\Gamma(d/2+1)$ = $d$ 维单位球体积。k-近邻版 $k(n)=k_{\mathrm{PRM}}\log n$，$k_{\mathrm{PRM}}>k_{\mathrm{PRM}}^*=e(1+1/d)$（只依赖 $d$，不依赖实例；$k=2e$ 通用）。
RRT$^*$ 用 $r(n)=\min\{\gamma_{\mathrm{RRT}^*}(\log n/n)^{1/d},\eta\}$，
\begin{equation}
  \gamma_{\mathrm{RRT}^*}>\big(2(1+\tfrac1d)\big)^{1/d}\Big(\frac{\mu(\mathcal{C}_{free})}{\zeta_d}\Big)^{1/d}.
  \label{eq:plan-gamma-rrt}
\end{equation}
k-近邻 RRT$^*$ 要 $k_{\mathrm{RRT}^*}>2^{d+1}e(1+1/d)$。
\end{definition}

\begin{pitfall}[概念误区：把 RRT$^*$ 与 PRM$^*$ 的常数当成一样]
注意 \cref{eq:plan-gamma-rrt} 的 RRT$^*$ 下界 $(2(1+1/d))^{1/d}$ 与 \cref{eq:plan-gamma-prm} 的 PRM$^*$/RRG 下界 $2(1+1/d)^{1/d}$ \textbf{指数位置不同}——RRT$^*$ 把因子 2 也放进 $1/d$ 次方内。这不是笔误：RRT$^*$ 是\emph{单查询增量树}，连接结构受“到达顺序”约束更强（见 \cref{sec:plan-aotheory} 的标记点过程），需要更大的半径常数才能保证最优。混用会导致半径取得偏小、丧失渐近最优。
\end{pitfall}

\begin{insight}
\textbf{为何 $(\log n/n)^{1/d}$ 是临界缩放？}\;
无障碍时 sPRM 是随机 $r$-disc 图。由 Penrose 的随机几何图连通性定理，$d$ 维随机 $r$-disc 图 $G^{disc}(n,r)$ 满足
\[
  \lim_{n\to\infty}\mathbb{P}\{G^{disc}(n,r)\text{ 连通}\}=\begin{cases}1,&\zeta_d r^d>\log n/n,\\0,&\zeta_d r^d<\log n/n.\end{cases}
\]
即连通的临界半径恰是 $r^d\sim\log n/(n\zeta_d)$——正是 $\gamma^*$ 公式里 $(\mu(\mathcal{C}_{free})/\zeta_d)^{1/d}$ 与 $(\log n/n)^{1/d}$ 的来源。半径再小（如 $n^{-1/d}$，缺 $\log$ 因子）则图落入\emph{亚临界渗流区}、碎成小簇够不到目标；半径过大则碰撞检测开销暴涨。$(\log n/n)^{1/d}$ 是兼顾连通/最优与计算开销的\emph{最小}缩放。这也呼应 \cref{sec:plan-prm} 练习里 $n$ 点弥散量级 $O((\log n/n)^{1/d})$。
\end{insight}

\subsection{概率完备性与渐近最优性}
\label{sec:plan-aotheory}

\begin{definition}[概率完备性]\label{def:plan-probcomplete}
称问题\textbf{鲁棒可行}，若存在 $\delta>0$ 使有\emph{强 $\delta$-clearance}（整条路径离障碍 $\ge\delta$，即 $\sigma(\tau)\in\mathrm{int}_\delta(\mathcal{C}_{free})$）的解。算法 ALG \textbf{概率完备}，若对任意鲁棒可行问题
\begin{equation}
  \liminf_{n\to\infty}\mathbb{P}\{\exists\mathbf{x}_{goal}\in V_n^{\mathrm{ALG}}\cap\mathcal{C}_{goal}\text{ 与 }\mathbf{x}_{init}\text{ 在 }G_n^{\mathrm{ALG}}\text{ 中连通}\}=1.
  \label{eq:plan-probcomplete}
\end{equation}
\end{definition}

\begin{theorem}[sPRM 与 RRT 概率完备（指数收敛）]\label{thm:plan-pc}
对鲁棒可行问题，存在仅依赖 $\mathcal{C}_{free},\mathcal{C}_{goal}$ 的常数 $a>0,n_0$ 使
\[
  \mathbb{P}\{\text{sPRM 在 }n\text{ 点时连通到目标}\}>1-e^{-an},\qquad
  \mathbb{P}\{V_n^{\mathrm{RRT}}\cap\mathcal{C}_{goal}\ne\varnothing\}>1-e^{-an},\quad\forall n>n_0.
\]
即失败概率随顶点数\emph{指数}趋零（前者 Kavraki 等 1998，后者 LaValle--Kuffner 2001）。RRG、RRT$^*$ 因构造上与 RRT 同顶点集且返回连通图，由此直接继承概率完备。
\end{theorem}

\begin{theorem}[1-最近邻 sPRM 不完备]\label{thm:plan-1nn}
$k=1$ 的 k-最近邻 sPRM \emph{不}概率完备，且 $\lim_{n\to\infty}\mathbb{P}\{\text{连通到目标}\}=0$。
\end{theorem}
\begin{proof}（梗概，Karaman--Frazzoli\cite{karaman2011sampling} 定理 17）
设 $L_n$ = 1-最近邻图总边长、$N_n$ = 连通分量数。由 Wade 2007，$L_n/n^{1-1/d}$ 与 $N_n/n$ 均方收敛到常数 $\gamma_L,\gamma_N$；故分量平均长度 $\widetilde L_n=L_n/N_n$ 经 Slutsky 定理满足 $n^{1/d}\widetilde L_n\xrightarrow{d}\gamma_L/\gamma_N$（常数）。含 $\mathbf{x}_{init}$ 的分量长度 $L_n'$ 与 $\widetilde L_n$ 同分布，故 $n^{1/d}L_n'$ 依概率收敛到常数，即 $L_n'\to0$ 依概率。取 $\epsilon<\inf_{\mathbf{x}\in\mathcal{C}_{goal}}\|\mathbf{x}-\mathbf{x}_{init}\|$：含起点分量缩到半径 $<\epsilon$，够不到目标，故连通到目标的概率 $\le\mathbb{P}\{L_n'>\epsilon\}\to0$。
\end{proof}

\noindent 类似地，半径 $r(n)=\gamma n^{-1/d}$（\emph{无} $\log$ 因子）的变半径 sPRM 也不概率完备（Penrose 渗流下界 + Borel--Cantelli）——这定量印证了 \cref{sec:plan-radius} 洞察里“缺 $\log$ 落入亚临界”的结论。

\begin{definition}[渐近最优性]\label{def:plan-ao}
设 $Y_n^{\mathrm{ALG}}$ = ALG 第 $n$ 步图中最优解代价，$c^*$ = 鲁棒最优解代价。ALG \textbf{渐近最优}，若对任意有有限代价鲁棒最优解的问题
\begin{equation}
  \mathbb{P}\Big\{\limsup_{n\to\infty}Y_n^{\mathrm{ALG}}=c^*\Big\}=1.
  \label{eq:plan-ao}
\end{equation}
因 $Y_n^{\mathrm{ALG}}\ge c^*$，这等价于 $Y_n^{\mathrm{ALG}}\to c^*$ 几乎必然。概率完备是渐近最优的\emph{必要}条件。
\end{definition}

\begin{theorem}[0-1 律]\label{thm:plan-01law}
若 ALG 终能找到可行解（$\limsup_n Y_n<\infty$），则 $\{\limsup_n Y_n=c^*\}$ 的概率非 0 即 1。
\end{theorem}
\begin{proof}
$\{\limsup_n Y_n=c^*\}=\{\limsup_{n\ge m}Y_n=c^*\}$ 对每个 $m$ 都成立，故它属于 $\{Y_n\}_{n\ge m}$ 生成的 $\sigma$-域 $\mathcal{F}'_m$ 对所有 $m$，即落在尾 $\sigma$-域 $\mathcal{T}=\bigcap_m\mathcal{F}'_m$。由 Kolmogorov 0-1 律，尾事件概率为 0 或 1。
\end{proof}
\noindent 即：采样式算法\emph{要么在几乎所有运行都收敛到最优，要么几乎都不}——不存在“一半运行收敛”的中间情形。

\begin{theorem}[RRT 不渐近最优——关键负面结果]\label{thm:plan-rrt-noopt}
标准 RRT \emph{不}渐近最优。因 RRT 每步图只增不减（$G_i^{\mathrm{RRT}}\subseteq G_{i+1}^{\mathrm{RRT}}$），极限 $\lim_n Y_n^{\mathrm{RRT}}=Y_\infty^{\mathrm{RRT}}$ 存在；结合 \cref{thm:plan-01law} 得它几乎必然\emph{严格大于} $c^*$：
\begin{equation}
  \mathbb{P}\Big\{\lim_{n\to\infty}Y_n^{\mathrm{RRT}}>c^*\Big\}=1.
  \label{eq:plan-rrt-suboptimal}
\end{equation}
\end{theorem}
\begin{proof}（梗概，Karaman--Frazzoli\cite{karaman2011sampling} 附录 B 引理 44--49\rebuilt）
取无障碍 $\mathcal{C}_{free}=[0,1]^d$、$\eta$ 足够大。称含根的第 $k$ 个孩子及其全部后代为树的“第 $k$ 分支”。\emph{必要条件}（引理 44）：渐近最优要求第 $k$ 分支对无穷多 $k$ 含某 $R$-球（心 $\mathbf{x}_{init}$，$0<R<\inf\|\mathbf{x}_{goal}-\mathbf{x}_{init}\|$）外的顶点——否则只有有限分支能伸出 $R$-球，最优代价被这有限分支的次优值卡住（用 \cref{def:plan-ao} 的零测假设排除“恰好最优”）。\emph{违反}：设 $L_k$ 为第 $k$ 分支的“最长首路径”长度。由引理 45（新样本到 $n$ 点最近邻期望距离 $n^{-1/d}$）与单调收敛，$\mathbb{E}[L_1]=\sum_{i\ge1}i^{-1/d}i^{-1}=\zeta(1+1/d)$（Riemann zeta，对所有 $d$ 有限），且 $\mathbb{E}[L_k]$ 单调降到 0；经 Markov 型不等式 $\mathbb{P}\{\sup_{\alpha\ge k}L_\alpha>\epsilon\}\le\mathbb{E}[L_k]/\epsilon\to0$。故第 $k$ 分支对无穷多 $k$ 伸出 $R$-球的概率为 0，违反必要条件，得 \cref{eq:plan-rrt-suboptimal}。
\end{proof}

\begin{insight}
\textbf{RRT 非最优的物理直觉}：RRT 一旦把某顶点连到树，\emph{永不重连}。早期建立的“歪”连接被冻结，后代只能继承；新分支的最长路径期望 $\to0$（$\zeta$ 函数尾部），意味着越往后的分支越“短”、够不到远处去改善已有路径。RRT$^*$ 的\emph{重接线}正是为打破这个冻结。这也解释了为何“多次重启 RRT 取最好”有效——相当于对随机变量 $Y_\infty^{\mathrm{RRT}}$ 多次采样。
\end{insight}

\begin{theorem}[PRM$^*$、RRG、RRT$^*$ 渐近最优]\label{thm:plan-ao-pos}
在 \cref{def:plan-radius} 的半径缩放（$\gamma>\gamma^*$ 或 $k>k^*$）下，sPRM、PRM$^*$、RRG、RRT$^*$ 均渐近最优——返回解代价以概率 1 收敛到 $c^*$。
\end{theorem}
\begin{proof}（球覆盖法骨架，Karaman--Frazzoli\cite{karaman2011sampling} 附录 C/G\rebuilt）
设鲁棒最优路径 $\sigma^*$（有弱 $\delta$-clearance）。(1) 构造强 $\delta_n$-clearance 路径列 $\sigma_n\to\sigma^*$（引理 50）。(2) 沿 $\sigma_n$ 铺一串\emph{覆盖球} $\{B_{n,m}\}$，每球半径 $\sim r(n)$、相邻球心相距 $\Theta(r(n))$，使任意相邻两球各取一点其连线段在 $\mathcal{C}_{free}$ 内、长度 $\le r(n)$。(3) 设 $A_n$ = “每球都含至少一个图顶点”，用二项大偏差 + $\gamma>\gamma^*$ 保证漏球概率可和，经 Borel--Cantelli 得 $\mathbb{P}(\liminf_n A_n)=1$。(4) $A_n$ 发生时相邻球顶点被边连成无碰撞路径，其代价 $\to c(\sigma_n)\to c(\sigma^*)=c^*$，由鲁棒最优性的代价连续性得 $\limsup_n Y_n=c^*$ 几乎必然。RRT$^*$ 的证明（附录 G）在球覆盖上多一层\emph{标记点过程}：给每个样本一个独立均匀的“到达顺序标记” $Y_i\in[0,1]$，连边须同时满足空间近邻（$\|\cdot\|\le r_n$）与时序 $Y_{i'}\le Y_i$；半径常数 $\gamma_{\mathrm{RRT}^*}$ 的更大下界恰好压过这个序统计的衰减率——这就是 \cref{eq:plan-gamma-rrt} 比 \cref{eq:plan-gamma-prm} 多一个因子的根本原因。
\end{proof}

\begin{note}
\textbf{RRT$^*$ 最优性证明的勘误（Solovey 等 2020）.}\;
Solovey、Janson、Schmerling、Frazzoli、Pavone\cite{solovey2020revisiting} 指出上述 RRT$^*$ 渐近最优性的原始证明存在一个逻辑缺口（覆盖球论证未充分处理样本的\emph{时序依赖}）。他们给出严格的替代证明，关键是把连接半径指数从 $1/d$ 调整为 $r(n)=\gamma'(\log n/n)^{1/(d+1)}$（计入决定样本排序的额外时间维度）。\emph{结论不变}——RRT$^*$ 仍渐近最优，但参数要求更严。正文用经典 $1/d$ 表述，此处并记最新更正以保严谨。
\end{note}

\subsection{复杂度与实验定论}
\label{sec:plan-complexity}

记 $M_n^{\mathrm{ALG}}$ = 第 $n$ 次迭代的碰撞检测调用数。

\begin{theorem}[每步碰撞检测复杂度]\label{thm:plan-complexity}
PRM、sPRM 为 $M_n\in\Omega(n)$（每步须对半径球内 $\propto n$ 个顶点试连）；k-最近 sPRM 为常数 $k$；RRT 为常数 1；而\textbf{PRM$^*$、RRG、RRT$^*$ 为 $O(\log n)$}。
\end{theorem}
\begin{proof}（引理 41/42 梗概）
sPRM：令 $\bar\mu=\inf_{\mathbf{x}}\mu(\mathcal{B}_{\mathbf{x},r}\cap\mathcal{C}_{free})>0$，第 $n$ 步对末样本半径 $r$ 球内节点试连，期望数 $\ge(\bar\mu/\mu(\mathcal{C}_{free}))n$，故 $\Omega(n)$。PRM$^*$：半径 $r_n=\gamma_{\mathrm{PRM}}(\log n/n)^{1/d}$，球内期望顶点数 $\approx\zeta_d\gamma_{\mathrm{PRM}}^d\log n$，即 $O(\log n)$。
\end{proof}

\begin{table}[H]\centering\small
\caption{采样算法对照：完备性、最优性、每步复杂度}
\label{tab:plan-algos}
\begin{tabular}{lccc}
\toprule
算法 & 结构 / 查询 & 渐近最优？ & $M_n$（每步碰撞检测） \\
\midrule
PRM & 森林 / 多查询 & 否（定理 29） & $\Omega(n)$ \\
sPRM & $r$-disc 图 / 多查询 & 是（定理 30） & $\Omega(n)$ \\
k-最近 sPRM & $k$-NN 图 / 多查询 & 否（定理 31） & $k$（常数） \\
RRT & 树 / 单查询 & \textbf{否（\cref{thm:plan-rrt-noopt}）} & $1$（常数） \\
RRT-Connect & 双树 / 单查询 & 否 & $O(1)$（贪心） \\
\textbf{PRM$^*$} & $r$-disc 图 / 多查询 & \textbf{是（\cref{thm:plan-ao-pos}）} & $O(\log n)$ \\
\textbf{RRG} & 图 / 单查询 & \textbf{是} & $O(\log n)$ \\
\textbf{RRT$^*$} & 树 / 单查询(anytime) & \textbf{是} & $O(\log n)$ \\
\bottomrule
\end{tabular}
\end{table}

\noindent\textbf{核心结论}：PRM$^*$/RRG/RRT$^*$ 以仅 $O(\log n)$ 的每步代价（仅是 RRT 常数的 $\log n$ 倍）达到渐近最优——“渐近最优几乎免费”，这是 Karaman--Frazzoli 的中心贡献。

\begin{insight}
\textbf{三个常被引用的实验数字}（Karaman--Frazzoli\cite{karaman2011sampling} §5 蒙特卡洛，各 20000 迭代 $\times$ 500 次）：
(1) \emph{无障碍}下 RRT 的解代价极限收敛到\textbf{最优的 $\sqrt2\approx1.414$ 倍}（确定性与随机设定均如此）；
(2) \emph{有障碍、双同伦类}场景下 RRT 平均约\textbf{最优的 $1.5$ 倍}（取决于碰巧落入哪个同伦类）；
(3) RRT$^*$ 方差 $\to0$、代价 $\to$ 最优，RRT 方差 $\to2.5$。
这定量印证 \cref{thm:plan-rrt-noopt}（RRT 非最优）与 \cref{thm:plan-ao-pos}（RRT$^*$ 最优）。RRT$^*$ 在含障碍场景里还会\emph{发现不同同伦类}的更短路径——这正是 \cref{sec:plan-topology} 同伦概念的实际意义。运行时间上 RRT$^*$/RRT 之比随迭代收敛到常数，符合 \cref{thm:plan-complexity}。
\end{insight}

\begin{practice}
\begin{enumerate}
  \item 解释 RRT$^*$ 的“重接线”如何打破 \cref{thm:plan-rrt-noopt} 证明中“分支首路径冻结”的机制：当一个更短的连接出现时，已有顶点的父如何更新？
  \item 给定 $d=2$、$\mu(\mathcal{C}_{free})=1$，算出 PRM$^*$ 的 $\gamma_{\mathrm{PRM}}^*$ 与 RRT$^*$ 的 $\gamma_{\mathrm{RRT}^*}^*$ 数值，验证后者更大（用 $\zeta_2=\pi$）。
  \item （开放思考）为什么“跑 5 个独立 RRT 取最短”常比单个 RRT 好，却仍\emph{不}渐近最优？用 \cref{eq:plan-rrt-suboptimal} 把 $Y_\infty^{\mathrm{RRT}}$ 当随机变量来论证。
\end{enumerate}
\end{practice}

% ════════════════════════════════════════════════════════════════
\section{轨迹优化简介}
\label{sec:plan-traj}

\paragraph{动机：搜索/采样给“哪条路”，优化给“多好的路”。}
搜索式（A$^*$/D$^*$）与采样式（PRM/RRT$^*$）解决\emph{可行性与全局结构}——在复杂 $\mathcal{C}_{free}$ 中找一条无碰撞路径。但它们的解常\emph{锯齿、不光滑、动力学不可行}（RRT 尤甚）。\textbf{轨迹优化}解决\emph{局部精化}：给定一条初始轨迹（可由 RRT$^*$ 给出，或简单直线初始化），通过最小化一个\emph{代价泛函}（含平滑度 + 障碍代价 + 约束）把它优化成光滑、低代价、动力学可行的轨迹。二者常\emph{级联}：采样式做全局、优化做局部。

\begin{definition}[轨迹优化的一般形式]\label{def:plan-trajopt}
求轨迹 $\xi:[0,1]\to\mathcal{C}$（或离散为 waypoint 序列 $\mathbf{q}_0,\dots,\mathbf{q}_{n+1}$，端点固定），最小化
\begin{equation}
  \min_{\xi}\ \mathcal{U}[\xi]=\underbrace{\mathcal{F}_{obs}[\xi]}_{\text{障碍代价}}+\lambda\underbrace{\mathcal{F}_{smooth}[\xi]}_{\text{平滑/动力学}}\quad\text{s.t. }\xi(0)=\mathbf{q}_I,\ \xi(1)=\mathbf{q}_G,\ (\text{动力学/关节限}),
  \label{eq:plan-trajopt}
\end{equation}
$\lambda$ = 平滑/避障权衡系数。平滑项常取导数平方积分 $\mathcal{F}_{smooth}=\tfrac12\int_0^1\|\tfrac{d}{dt}\xi(t)\|^2dt$（离散为 $\tfrac12\xi^\top\mathbf{A}\xi+\xi^\top\mathbf{b}+c$，$\mathbf{A}=\mathbf{K}^\top\mathbf{K}$ 由有限差分矩阵构成、对称正定，度量轨迹总加速度）；障碍项由到障碍的距离场构造。
\end{definition}

\subsection{CHOMP：协变梯度下降}
\label{sec:plan-chomp}

CHOMP（Covariant Hamiltonian Optimization for Motion Planning，Zucker 等 2013\cite{zucker2013chomp}）是基于梯度的轨迹优化代表。其障碍泛函在\emph{工作空间}做弧长参数化：设 $\mathcal{B}\subset\mathbb{R}^3$ 为机器人体表点集，$\mathbf{x}(\mathbf{q},u)$ 为正运动学（构型 + 体点 $\mapsto$ 工作空间点），$c:\mathbb{R}^3\to\mathbb{R}$ 为工作空间代价场（障碍内及附近惩罚大，由到障碍的欧氏距离变换 EDT 定义），则
\begin{equation}
  \mathcal{F}_{obs}[\xi]=\int_0^1\!\!\int_{\mathcal{B}}c\big(\mathbf{x}(\xi(t),u)\big)\,\Big\|\tfrac{d}{dt}\mathbf{x}(\xi(t),u)\Big\|\,du\,dt.
  \label{eq:plan-chomp-obs}
\end{equation}
代价 $c$ 乘以体点工作空间速度的范数，把线积分变为\emph{弧长参数化}，使障碍目标对重新计时不变（同一路径不同速度走，$\mathcal{F}_{obs}$ 不变）。

\begin{derivation}[CHOMP 的泛函梯度（Euler--Lagrange）]\label{der:plan-chomp}
设 $\mathcal{U}[\xi]=\int v(\xi(t),\xi'(t))\,dt$，扰动 $\bar\xi=\xi+\epsilon\phi$（$\phi(0)=\phi(1)=0$，端点固定），视 $\mathcal{U}[\bar\xi]=f(\epsilon)$。对第二项分部积分：
\[
  \frac{df}{d\epsilon}=\int\Big(\frac{\partial v}{\partial\bar\xi}\cdot\phi+\frac{\partial v}{\partial\bar\xi'}\cdot\phi'\Big)dt
  =\int\Big(\frac{\partial v}{\partial\bar\xi}-\frac{d}{dt}\frac{\partial v}{\partial\bar\xi'}\Big)\cdot\phi\,dt,
\]
末步用 $\phi(0)=\phi(1)=0$ 消去边界项。令 $\epsilon=0$，使该方向导数最大（在有限欧氏范数下）的 $\phi$ 正比于括号项，故\textbf{泛函梯度}
\begin{equation}
  \bar\nabla\mathcal{U}[\xi]=\frac{\partial v}{\partial\xi}-\frac{d}{dt}\frac{\partial v}{\partial\xi'},
  \label{eq:plan-funcgrad}
\end{equation}
这正是\emph{Euler--Lagrange 方程}（$\bar\nabla\mathcal{U}=0$ 即临界点）。平滑项（速度平方）梯度为 $\bar\nabla\mathcal{F}_{smooth}=-\tfrac{d^2}{dt^2}\xi$。
\end{derivation}

\begin{theorem}[CHOMP 协变更新规则]\label{thm:plan-chomp-update}
普通梯度下降 $\xi_{i+1}=\xi_i-\eta\bar\nabla\mathcal{U}$ 依赖轨迹\emph{表示}（其泛函梯度源于欧氏范数）。CHOMP 改用只依赖轨迹动态量的范数 $\|\xi\|_A^2=\langle\xi,\mathbf{A}\xi\rangle=\int(D\xi)^2dt$（$\mathbf{A}=D^\dagger D$）。把更新写成“使 $\mathcal{U}$ 下降最多、同时在度量 $\mathbf{A}$ 下与 $\xi_i$ 接近”的优化
\[
  \xi_{i+1}=\arg\min_\xi\ \mathcal{U}[\xi_i]+(\xi-\xi_i)^\top\bar\nabla\mathcal{U}[\xi_i]+\frac\eta2\|\xi-\xi_i\|_A^2,
\]
对 $\xi$ 求导置零得\textbf{协变更新规则}
\begin{equation}
  \xi_{i+1}=\xi_i-\frac1\eta\,\mathbf{A}^{-1}\bar\nabla\mathcal{U}[\xi_i].
  \label{eq:plan-chomp-update}
\end{equation}
\end{theorem}

\begin{insight}
\cref{eq:plan-chomp-update} 的 $\mathbf{A}^{-1}\bar\nabla\mathcal{U}$ 与本书 \cref{ch:nlopt} 的高斯牛顿/自然梯度\emph{同框架}：$\mathbf{A}=\mathbf{K}^\top\mathbf{K}$ 即平滑先验的信息矩阵（类比 $\mathbf{H}=\mathbf{J}^\top\mathbf{J}$），$\mathbf{A}^{-1}\nabla$ 即一步“自然梯度”。直觉上，普通梯度对“单点障碍冲击”只调那一点的 waypoint，而 $\mathbf{A}^{-1}$ 把这个冲击\emph{协变地分摊到轨迹一大段}，使调整后仍光滑——预期比各向同性梯度法快约 $O(n)$ 倍收敛。若轨迹在 $\SE(3)$ 上优化，平滑项与梯度须用 \cref{ch:lie} 的右雅可比 $\mathbf{J}_r$。
\end{insight}

\paragraph{代表方法与级联流水线。}\rebuilt
\textbf{CHOMP} 用协变梯度下降，障碍项基于工作空间距离场；无完备性（可能以仍有碰撞的轨迹终止），故常配 Hamiltonian Monte Carlo 缓解局部极小并获概率完备。\textbf{STOMP}（Kalakrishnan 等 2011）\emph{无梯度}，采样一批噪声轨迹按代价加权更新，适合不可微代价/硬约束。\textbf{TrajOpt}（Schulman 等 2013）把轨迹优化写成序列凸规划（SQP），碰撞约束用连续碰撞凸近似，处理硬约束更稳健。\textbf{KOMO / 因子图轨迹优化}把轨迹优化表为非线性最小二乘 $\min\sum_k\|f_k(\xi)\|^2$，用 GN/LM 求解，与本书 SLAM 后端天然衔接。实际系统常\textbf{级联}：A$^*$/RRT$^*$ 给全局初值 $\to$ CHOMP/TrajOpt 局部平滑成可执行轨迹。

\begin{table}[H]\centering\small
\caption{采样式 vs 轨迹优化}
\label{tab:plan-traj-cmp}
\begin{tabular}{lll}
\toprule
维度 & 采样式（RRT$^*$/PRM$^*$） & 轨迹优化（CHOMP） \\
\midrule
完备性 & 概率完备 / 渐近最优 & 局部法，一般不完备（+HMC 才概率完备） \\
初值 & 不需可行初值（从根长树） & 需初始轨迹（常用直线，可不可行） \\
解性质 & 全局、探索多同伦类；常 jagged 需后平滑 & 局部最优、天然光滑；不跨同伦类 \\
典型用法 & 找全局粗解 & 平滑/精化其他方法的结果 \\
\bottomrule
\end{tabular}
\end{table}

\begin{practice}
\begin{enumerate}
  \item 从 \cref{der:plan-chomp} 出发，独立推出平滑项 $\mathcal{F}_{smooth}=\tfrac12\int\|\xi'\|^2dt$ 的泛函梯度 $-\xi''$（在草稿纸上完成分部积分）。
  \item 解释 \cref{eq:plan-chomp-update} 中 $\mathbf{A}^{-1}$ 的作用：为什么用 $\mathbf{A}$（总加速度度量）而非单位阵 $\mathbf{I}$ 会让更新“天然光滑”？
  \item （跨章综合）把一段 $\SE(3)$ 上的位姿轨迹平滑写成最小二乘问题，指出平滑残差对位姿的雅可比应使用 \cref{ch:lie} 的哪个量（右扰动 $\mathbf{J}_r$），并说明它如何接入 \cref{ch:nlopt} 的高斯牛顿。
\end{enumerate}
\end{practice}

% ════════════════════════════════════════════════════════════════
\section*{本章常见误解汇总}
\begin{table}[H]\centering\small
\begin{tabular}{p{0.46\textwidth} p{0.46\textwidth}}
\toprule
\textbf{常见误解} & \textbf{正确理解} \\
\midrule
$\mathcal{C}_{free}$ 是闭集，可直接谈“最短路” & $\mathcal{C}_{free}$ 开（$\mathcal{C}_{obs}$ 闭），最优化须用 $\mathrm{cl}(\mathcal{C}_{free})$（\cref{def:plan-cobs}）\\
$\mathbf{R}^\top\mathbf{R}=\mathbf{I}$ 即旋转 & 还需 $\det=+1$，否则含反射（$O(3)=S^3\sqcup S^3$）\\
欧拉角能干净描述 $\SO(3)$ & 万向锁/拓扑破坏；须用四元数，$\SO(3)\cong\mathbb{RP}^3$（\cref{eq:plan-so3rp3}）\\
高维直接用细栅格 + A$^*$ & 维数灾难 $O(1/h^d)$；六自由度即不可行（\cref{thm:plan-pspace}）\\
A$^*$ 与 best-first 都用 $h$ 故差不多 & A$^*$ 计入 $g$ 故最优；best-first 只用 $h$、贪心不保最优 \\
可采纳启发在图搜索就够 & 图搜索保最优需\emph{一致}（\cref{thm:plan-cons-mono}）\\
平方欧氏距离是更精的启发 & 量纲不匹配致高估，破坏可采纳（\cref{sec:plan-grid}）\\
RRT 跑久了会收敛到最优 & RRT 几乎必然收敛到次优值（\cref{thm:plan-rrt-noopt}）\\
RRT$^*$ 与 PRM$^*$ 的 $\gamma^*$ 公式相同 & 指数位置不同：RRT$^*$ 为 $(2(1+1/d))^{1/d}$（\cref{eq:plan-gamma-rrt}）\\
半径取 $\gamma n^{-1/d}$ 也行（省 $\log$） & 落入亚临界渗流、不完备；须 $(\log n/n)^{1/d}$ \\
轨迹优化能独立解全局规划 & 局部法不完备；常需采样式给初值再级联 \\
\bottomrule
\end{tabular}
\end{table}

% ════════════════════════════════════════════════════════════════
\section*{本章小结}

\begin{table}[H]\centering\small
\caption{符号表}
\begin{tabular}{lll}
\toprule
符号 & 含义 & 首次出现 \\
\midrule
$\mathcal{C},\mathcal{C}_{free},\mathcal{C}_{obs}$ & 构型空间 / 自由空间 / 障碍空间 & \cref{eq:plan-cobsfree} \\
$\mathbf{q}_I,\mathbf{q}_G,\mathcal{C}_{goal}$ & 起点 / 目标 / 目标集 & \cref{def:plan-piano} \\
$\tau,\sigma,\xi$ & 路径 / 路径（采样源）/ 轨迹（优化源） & \cref{def:plan-piano} \\
$\mathbb{RP}^3$ & 三维实射影空间（$\cong\SO(3)$） & \cref{eq:plan-rpn,eq:plan-so3rp3} \\
$g,h,f$ & cost-to-come / 启发 / 评价 $f=g+h$ & \cref{eq:plan-fgh} \\
$h^*$ & 真实最优 cost-to-go & \cref{eq:plan-admissible} \\
$G^*(\mathbf{x})$ & 最优 cost-to-go（导航函数） & \cref{thm:plan-bellman} \\
$\eta$ & Steer 步长上界 & \cref{sec:plan-primitives} \\
$r(n),\gamma,k$ & 连接半径 / 半径常数 / 近邻数 & \cref{def:plan-radius} \\
$\zeta_d,\mu$ & $d$ 维单位球体积 / Lebesgue 测度 & \cref{eq:plan-gamma-prm} \\
$Y_n^{\mathrm{ALG}}$ & 第 $n$ 步图中最优解代价 & \cref{def:plan-ao} \\
$c^*$ & 鲁棒最优解代价 & \cref{def:plan-ao} \\
$\mathbf{A}=\mathbf{K}^\top\mathbf{K}$ & 轨迹平滑度量（总加速度） & \cref{def:plan-trajopt} \\
$\bar\nabla\mathcal{U}$ & 泛函梯度 & \cref{eq:plan-funcgrad} \\
\bottomrule
\end{tabular}
\end{table}

\begin{table}[H]\centering\small
\caption{定理速查表}
\begin{tabular}{lll}
\toprule
定理 / 公式 & 一句话说明 & 对应 \\
\midrule
钢琴搬运工 PSPACE-hard & 精确规划随维数爆炸，逼出采样法 & \cref{thm:plan-pspace} \\
Bellman 递推 & cost-to-go 的最优性原理递推 & \cref{thm:plan-bellman} \\
Dijkstra 出队即最优 & 队首代价已最优（非负权） & \cref{thm:plan-dijkstra} \\
A$^*$ 可采纳$\Rightarrow$最优 & $h\le h^*$ 保证返回最短路 & \cref{thm:plan-astar-opt} \\
一致$\Rightarrow$可采纳 + 单调 & 三角不等式保无重复展开 & \cref{thm:plan-cons-adm,thm:plan-cons-mono} \\
临界半径 $(\log n/n)^{1/d}$ & 连通/最优的最小连接缩放 & \cref{def:plan-radius} \\
RRT 不渐近最优 & 解几乎必然收敛到次优值 & \cref{thm:plan-rrt-noopt} \\
PRM$^*$/RRG/RRT$^*$ 渐近最优 & $O(\log n)$ 代价换最优 & \cref{thm:plan-ao-pos} \\
CHOMP 协变更新 & $\xi\leftarrow\xi-\tfrac1\eta\mathbf{A}^{-1}\bar\nabla\mathcal{U}$ & \cref{thm:plan-chomp-update} \\
\bottomrule
\end{tabular}
\end{table}

\begin{table}[H]\centering\small
\caption{知识点总表}
\begin{tabular}{clll}
\toprule
\# & 知识点 & 核心要点 & 难度 \\
\midrule
1 & 构型空间 & 刚体避障 $\to$ 点找路；$\SO(3)\cong\mathbb{RP}^3$ & \(\star\star\) \\
2 & 钢琴搬运工问题 & PSPACE-hard，显式 $\mathcal{C}_{obs}$ 不可行 & \(\star\star\) \\
3 & 图搜索框架 & 仅 $Q$ 排序不同；BFS/DFS/Dijkstra/A$^*$ & \(\star\star\) \\
4 & 动态规划 & Bellman 递推、导航函数 & \(\star\star\star\) \\
5 & A$^*$ 与启发 & 可采纳 $\Rightarrow$ 最优、一致 $\Rightarrow$ 无重复展开 & \(\star\star\star\) \\
6 & 栅格启发 & 曼哈顿/对角/欧氏、平方陷阱 & \(\star\star\) \\
7 & RRT / RRT-Connect & Voronoi 偏置、双树贪心连接 & \(\star\star\star\) \\
8 & PRM & 多查询路标、$r$-disc 图、union-find & \(\star\star\star\) \\
9 & RRT$^*$/RRG/PRM$^*$ & 选最优父 + 重接线、临界半径 & \(\star\star\star\star\) \\
10 & 渐近最优理论 & 概率完备/渐近最优、RRT 非最优证明 & \(\star\star\star\star\star\) \\
11 & 轨迹优化 & CHOMP 协变梯度、级联流水线 & \(\star\star\star\star\) \\
\bottomrule
\end{tabular}
\end{table}

% ════════════════════════════════════════════════════════════════
\section*{延伸阅读}
\begin{itemize}
  \item \textbf{规划全景与构型空间}：LaValle《Planning Algorithms》\cite{lavalle2006planning}（官方全文 lavalle.pl/planning）——拓扑、流形、$\SO(3)/\SE(3)$、离散搜索、采样式规划的奠基教材，本章 C-space、图搜索、RRT 的主要出处。
  \item \textbf{渐近最优采样规划}：Karaman \& Frazzoli《Sampling-based Algorithms for Optimal Motion Planning》\cite{karaman2011sampling}（arXiv:1105.1186）——PRM$^*$/RRG/RRT$^*$ 算法、概率完备/渐近最优定理与完整证明、临界半径常数的权威来源。
  \item \textbf{A$^*$ 原典}：Hart, Nilsson, Raphael《A Formal Basis for the Heuristic Determination of Minimum Cost Paths》\cite{hart1968astar}——$g+h$、可采纳性、一致性的起源；栅格启发实现见 Patel 的斯坦福教程\cite{patel2018heuristics}。
  \item \textbf{单查询双树}：Kuffner \& LaValle《RRT-Connect》\cite{kuffner2000rrtconnect}（ICRA 2000）——双树 + 贪心 CONNECT 的工程首选规划器、含自包含完备性证明。
  \item \textbf{动态重规划}：Stentz《D$^*$》\cite{stentz1994optimal}（ICRA 1994）——部分已知/动态环境的增量最优搜索；现代实践用其后继 D$^*$ Lite。
  \item \textbf{RRT$^*$ 最优性勘误}：Solovey 等《Revisiting the Asymptotic Optimality of RRT$^*$》\cite{solovey2020revisiting}（ICRA 2020）——指出原证明缺口、给出 $(\log n/n)^{1/(d+1)}$ 修正半径，近年最重要的理论更正。
  \item \textbf{轨迹优化}：Zucker 等《CHOMP》\cite{zucker2013chomp}（IJRR 2013）——协变梯度轨迹优化的目标泛函与更新规则；STOMP、TrajOpt 为同类代表。
\end{itemize}

\section*{本章与后续章节的关系}
\begin{table}[H]\centering\small
\begin{tabular}{lll}
\toprule
后续章节 & 关系 & 本章铺垫 \\
\midrule
\cref{ch:control} 控制导论 & 轨迹跟踪 / 动力学可行轨迹的执行 & 轨迹优化、级联流水线 \\
\cref{ch:lie} 李群（前置） & $\SE(3)$ 上规划的距离 / 插值 / 梯度 & C-space 流形、右扰动雅可比 \\
\cref{ch:nlopt} 非线性优化（前置） & 轨迹优化 = 最小二乘 + GN/LM & CHOMP 的 $\mathbf{A}^{-1}\nabla$ 与因子图轨迹优化 \\
\bottomrule
\end{tabular}
\end{table}

\section*{故障排查手册}
\begin{enumerate}
  \item \textbf{症状}：A$^*$ 给出的不是最短路。\textbf{排查}：检查启发是否\emph{高估}（如误用平方欧氏距离，\cref{sec:plan-grid}）；图搜索时检查 $h$ 是否\emph{一致}（仅可采纳不够，须 \cref{thm:plan-cons-mono}）；确认边代价非负。
  \item \textbf{症状}：RRT 在狭窄通道里长时间找不到解。\textbf{原因}：单树连目标靠目标偏置碰运气、bug trap 难穿。\textbf{对策}：改用 RRT-Connect 双树相向生长（\cref{sec:plan-rrtconnect}）；适度调高目标偏置概率。
  \item \textbf{症状}：RRT$^*$ 跑很久仍不收敛到最优。\textbf{排查}：连接半径常数 $\gamma_{\mathrm{RRT}^*}$ 是否取得过小（误用 PRM$^*$ 公式，\cref{eq:plan-gamma-rrt}）；是否漏写“选最优父 + 重接线”任一步；严谨场景按 Solà-Solovey 用 $1/(d+1)$ 指数（\cref{sec:plan-aotheory} note）。
  \item \textbf{症状}：PRM/sPRM 在高维下极慢。\textbf{原因}：每步碰撞检测 $\Omega(n)$（\cref{thm:plan-complexity}）。\textbf{对策}：改用 PRM$^*$ 的 $r(n)=\gamma(\log n/n)^{1/d}$ 把每步降到 $O(\log n)$；用 Kd-树加速最近邻。
  \item \textbf{症状}：CHOMP 终止时轨迹仍穿过障碍。\textbf{原因}：CHOMP 无完备性，陷入高代价局部极小。\textbf{对策}：用更好的初值（RRT$^*$ 全局解）做级联；加 Hamiltonian Monte Carlo 扰动；检查工作空间距离场是否在障碍内给出正确的排斥梯度。
\end{enumerate}
